% !TeX root = ../../Book1.tex
\chapter{分布与弱导数}
\label{b1:ch:20}
\BookChapterNumber{b1:ch:20}
% 本章摘要（不计入编号）
\begin{chapterabstract}
分部积分把求导转移到测试函数上，使不可微函数乃至点质量也能拥有导数。本章构造紧支集光滑测试空间，定义分布及其基本操作，证明分布的局部光滑化，再从中识别能够由局部可积函数表示的弱导数。尖角说明少量经典不可微点并不妨碍弱导数存在；跳跃函数与 Cantor 函数则说明，仅有经典导数几乎处处存在仍不足以保证存在局部可积弱导数。
\end{chapterabstract}

% 本章目标（不计入编号）
\begin{learninggoals}
\begin{itemize}
\item\label{b1:item:20-goal-tests}能使用光滑截断与单位分解，解释测试函数收敛为何同时要求共同紧支集和各阶导数控制。
\item\label{b1:item:20-goal-distributions}能从连续线性泛函的定义计算分布导数、限制、支集和光滑乘法，区分点质量与普通函数。
\item\label{b1:item:20-goal-smoothing}能说明分布与光滑核卷积的定义域、微分公式和收敛意义，避免把局部光滑化误作任意延拓。
\item\label{b1:item:20-goal-weak}能检验弱导数的存在性与唯一性，连接经典微分、局部 Lipschitz 函数和一维绝对连续代表。
\end{itemize}
\end{learninggoals}

从\secref{b1:sec:2001}到\secref{b1:sec:2004}，测试空间、其连续对偶以及对偶中仍由函数表示的部分依次出现。测试空间的局部凸拓扑只展开到本章需要的有限阶估计；读者第一次阅读时可以先掌握共同紧支集的收敛判据，再回看其拓扑构造。主要前置工具是\chapref{b1:ch:13}的卷积与 \(L^p\) 平移连续性、\chapref{b1:ch:14}的 Lebesgue 点定理，以及\chapref{b1:ch:16}的 Rademacher 定理和一维绝对连续函数理论。

Hunter 与 Nachtergaele 的《Applied Analysis》\cite[Chapter 11 and Section 12.9]{Hunter2001Applied}适合对照对偶配对、导数和典型例子，其中 Chapter 11 主要面向温和分布，本章则在一般开集上使用紧支集测试函数，不施加无穷远增长条件。Schwartz 的《广义函数论》\cite[第一至三章]{Schwartz2010DistributionsChinese}把测试空间、局部化和导数系统地联系起来，可用作进一步阅读。Sobolev 空间的应用路线可接 Evans 的《Partial Differential Equations》\cite[Chapter 5]{Evans2010Partial}；本章先完成这条路线所需的弱导数语言。

% !TeX root = ../../Book1.tex
\section{测试函数}
\label{b1:sec:2001}

对光滑函数作分部积分，可以把一个函数的导数移到另一个函数上。若第二个函数在边界附近为零，边界项也随之消失。于是，即使第一个函数不可微，积分的另一边仍可能有意义。为了沿着这条思路推广导数，先要选择一类既能任意次求导、又能把计算限制在开集内部的函数。

\ProseParagraphBreak

本节固定非空开集 \(\Omega\subset\mathbb R^d\)，标量域为 \(\mathbb K=\mathbb R\) 或 \(\mathbb C\)。记号 \(K\Subset\Omega\) 表示 \(K\) 是包含于 \(\Omega\) 的紧集。本节的拓扑构造只使用半范数、紧集穷竭和光滑截断，不预借分布或 Sobolev 空间的性质。
% 来源边界：Hunter2001Applied第11章主要使用Schwartz空间；作者公开稿印刷294页明确区分D'与S'。
% 本节DK完备性、单位分解和D的局部凸拓扑由下文完整构造，不把该书温和分布版本冒称为这里的证明。

\subsection{紧支集光滑函数}
\label{b1:subsec:200101}

对多重指标 \(\alpha=(\alpha_1,\ldots,\alpha_d)\in\mathbb N_0^d\)，约定
\[
 |\alpha|\coloneqq\sum_{i=1}^d\alpha_i,\quad
 \alpha!\coloneqq\prod_{i=1}^d\alpha_i!,\quad
 D^\alpha\coloneqq\partial_1^{\alpha_1}\cdots\partial_d^{\alpha_d}.
\]
不等式 \(\beta\leq\alpha\) 指每个坐标分别满足不等式，且
\(\binom\alpha\beta=\alpha!/[\beta!(\alpha-\beta)!]\)。这里 \(D^0\) 为恒等操作。

\begin{definition}\label{b1:def:test-function}
设 \(\varphi:\Omega\rightarrow\mathbb K\) 光滑。如果其支集
\[
 \operatorname{supp}\varphi
 \coloneqq\overline{\{\,x\in\Omega:\varphi(x)\ne0\,\}}^{\,\Omega}
\]
为 \(\Omega\) 的紧子集，则称 \(\varphi\) 为 \(\Omega\) 上的\term[ceshi-hanshu]{测试函数}[test function]。全体测试函数构成的线性空间记作
\[
 \mathcal D(\Omega)\coloneqq C_c^\infty(\Omega).
\]
\end{definition}
\symidx{\(\mathcal D(\Omega)\)：紧支集光滑测试函数空间}

测试函数也有“试验函数”这一译名；《分析学》第2版中译本第6.2节采用后者\footnote{该译名及节次可核对\href{https://academic.hep.com.cn/br/CN/chapter/978-7-04-017381-9/chapter06}{高等教育出版社公布的第六章目录}。这里仅借目录说明名称，不以目录作为下文拓扑构造的依据。}\termidx[shiyan-hanshu]{试验函数}。

\ProseParagraphBreak

支集的紧性和它位于开集内部同样重要。例如 \(\Omega=(0,1)\) 上的常函数 \(1\) 不是测试函数，尽管它处处光滑且有界。测试函数在靠近 \(\partial\Omega\) 的某个区域内为零，因此在 \(\Omega\) 外补零后仍是 \(\mathbb R^d\) 上的光滑函数。以后对测试函数取全空间上确界或积分时，均使用这个零延拓。

\ProseParagraphBreak

这种函数并非只有零函数。令
\[
 \vartheta(t)=
 \begin{cases}e^{-1/t},&t>0,\\0,&t\leq0.\end{cases}
\]
在正半轴上，每阶导数均是 \(e^{-1/t}\) 与 \(1/t\) 的某个多项式的乘积。由 \(s^Ne^{-s}\to0\)，这些导数在 \(t\downarrow0\) 时都趋于零；逐阶检查零点的差商，得到 \(\vartheta\in C^\infty(\mathbb R)\)。于是 \(x\mapsto\vartheta(1-|x|^2)\) 是支集为闭单位球、在开单位球内严格为正的光滑函数。平移与缩放便能在 \(\Omega\) 内任一球上放置这样的函数。

\subsection{基本操作}
\label{b1:subsec:200102}

若 \(\varphi\in\mathcal D(\Omega)\)，则 \(D^\alpha\varphi\in\mathcal D(\Omega)\)，且微分不会扩大支集。若 \(a\in C^\infty(\Omega)\)，则 \(a\varphi\) 也是测试函数；即使 \(a\) 无界，其在 \(\operatorname{supp}\varphi\) 附近的每个固定阶导数仍有界。Leibniz 公式为
\begin{equation}\label{b1:eq:test-leibniz}
 D^\alpha(a\varphi)
 =\sum_{\beta\leq\alpha}\binom\alpha\beta
 D^\beta a\,D^{\alpha-\beta}\varphi.
\end{equation}
这由一阶乘积法则对 \(|\alpha|\) 归纳得到，二项系数记录各次微分落在两个因子上的分配方式。

\ProseParagraphBreak

对 \(K\Subset\Omega\) 记
\[
 \mathcal D_K=\{\,\varphi\in\mathcal D(\Omega):\operatorname{supp}\varphi\subset K\,\},
 \quad p_{K,m}(\varphi)=\max_{|\alpha|\leq m}\sup_{x\in K}|D^\alpha\varphi(x)|.
\]
上确界可等价地取遍 \(\Omega\)，空紧集上的半范数取零。因为导数在 \(K\) 外为零，\eqref{b1:eq:test-leibniz}给出
\[
 p_{K,m}(a\varphi)\leq C_{a,K,m}\cdot p_{K,m}(\varphi),
 \quad p_{K,m}(D^\alpha\varphi)\leq p_{K,m+|\alpha|}(\varphi).
\]
常数只涉及 \(a\) 在 \(K\) 上不超过 \(m\) 阶的导数。

\ProseParagraphBreak

再记 \((\tau_h\varphi)(x)=\varphi(x-h)\)。当 \(h\to0\) 且充分小时，这些平移的支集落在 \(\Omega\) 的同一个紧子集中。各阶导数的一致连续性给出 \(D^\alpha(\tau_h\varphi-\varphi)\to0\) 一致成立。沿坐标方向还可用一维积分公式写出
\[
 \frac{\varphi(x-he_i)-\varphi(x)}h
 =-\int_0^1\partial_i\varphi(x-the_i)\,\mathrm dt.
\]
对任意阶导数重复这一等式，便知差商以相同方式趋于 \(-\partial_i\varphi\)。这比只知道每个点的差商收敛更强：后面连续泛函正需要这种各阶一致控制。

\begin{lemma}\label{b1:lem:smooth-cutoff}
设 \(K\Subset\Omega\)，\(U\subset\Omega\) 为包含 \(K\) 的开集。存在 \(\chi\in\mathcal D(U)\)，满足 \(0\leq\chi\leq1\)，且 \(\chi\) 在 \(K\) 的一个邻域内恒等于 \(1\)。
\end{lemma}
\begin{proof}
由上面的 \(\vartheta\) 构造光滑阶跃
\[
 b(t)=\frac{\vartheta(2-t)}{\vartheta(2-t)+\vartheta(t-1)}.
\]
分母处处为正，且 \(b=1\) 于 \((-\infty,1]\)，\(b=0\) 于 \([2,\infty)\)。选有限个球 \(B(x_j,r_j)\) 覆盖 \(K\)，使 \(\overline B(x_j,\sqrt2r_j)\subset U\)。令 \(b_j(x)=b(|x-x_j|^2/r_j^2)\)，并取
\[
 \chi=1-\prod_j(1-b_j).
\]
每个 \(b_j\) 在其小球上等于 \(1\)，在大球外为零。因此 \(\chi\) 具有所需值域、支集和邻域恒等性质。空紧集情形取 \(\chi=0\)。
\end{proof}

光滑截断把一个局部对象变成紧支集对象，但不会在截断恒等于 \(1\) 的区域内改变它。这个局部性质比截断函数的具体公式更常用。

\subsection{单位分解}
\label{b1:subsec:200103}

\begin{definition}\label{b1:def:smooth-partition-unity}
设 \((U_\lambda)_{\lambda\in\Lambda}\) 为 \(\Omega\) 的开覆盖。如果一族非负光滑函数 \((\psi_j)\) 的支集在 \(\Omega\) 中局部有限，每个支集包含于某个 \(U_{\lambda(j)}\)，并且
\[
 \sum_j\psi_j(x)=1\quad(x\in\Omega),
\]
则称它为从属于该覆盖的\term[guanghua-danwei-fenjie]{光滑单位分解}[smooth partition of unity]。局部有限指每个点有一个邻域，只与有限个支集相交。
\end{definition}

\begin{theorem}\label{b1:thm:smooth-partition-unity}
每个欧氏开集的开覆盖都有一族可数、从属于该覆盖的光滑单位分解，而且每个分解函数均有紧支集。
\end{theorem}
\begin{proof}
取紧集穷竭
\[
 K_j=\{\,x\in\Omega:|x|\leq j,\ \operatorname{dist}(x,\mathbb R^d\setminus\Omega)\geq1/j\,\},
 \quad j\geq1,
\]
空补集的距离约定为 \(+\infty\)，并令 \(K_0=K_{-1}=\varnothing\)。这些紧集满足 \(K_j\subset\operatorname{int}K_{j+1}\)，且其内部之并为 \(\Omega\)。紧集
\(A_j=K_j\setminus\operatorname{int}K_{j-1}\) 包含在开集 \(\operatorname{int}K_{j+1}\setminus K_{j-2}\) 内。

对每个 \(A_j\)，选有限个小球覆盖它，使各小球的稍大闭球仍包含在 \(\operatorname{int}K_{j+1}\setminus K_{j-2}\) 与某个 \(U_\lambda\) 的交中。由\cref{b1:lem:smooth-cutoff}，可取在小球闭包上等于 \(1\)、支集在相应大球内的非负光滑函数 \(\eta_{j,l}\)。这族函数的支集局部有限：给定点先取其邻域包含在某个 \(\operatorname{int}K_N\) 中，所有 \(j\geq N+2\) 的支集均避开该邻域，其余层只有有限个函数。

各层小球覆盖 \(\Omega\)，故局部有限和 \(s=\sum_{j,l}\eta_{j,l}\) 光滑且处处为正。令 \(\psi_{j,l}=\eta_{j,l}/s\)，再把双指标排成一个序列。除以正光滑函数不扩大支集，所需各性质随即成立。
\end{proof}

\begin{corollary}\label{b1:cor:precompact-smooth-partition}
若 \(\Omega\subseteq\mathbb R^d\) 为非空开集，则存在可数的光滑单位分解 \((\psi_i)\) 及局部有限开覆盖 \((V_i)\)，使
\[
 \psi_i\in C_c^\infty(V_i),\qquad V_i\Subset\Omega.
\]
特别地，每个紧集只与有限多个 \(V_i\) 相交。
\end{corollary}
\begin{proof}
在上一证明中保留承载各个 \(\eta_{j,l}\) 的稍大开球，并与相应的 \(\psi_{j,l}\) 一同枚举。这些开球的闭包均包含在 \(\Omega\) 中；每层只有有限个，而足够远的层避开任意给定紧集，故它们构成局部有限开覆盖。归一化不扩大支集，所以相应的 \(\psi_i\) 属于 \(C_c^\infty(V_i)\)。
\end{proof}

若 \(\varphi\) 是测试函数，局部有限性和 \(\operatorname{supp}\varphi\) 的紧性使
\[
 \varphi=\sum_j\psi_j\varphi
\]
实际上只有有限个非零项。于是可以先在每个覆盖开集上验证一个关于测试函数的线性恒等式，再把有限项相加恢复全局恒等式。局部有限的要求不能换成“可数”二字；一般可数和既未必光滑，也未必允许逐项操作。

\subsection{测试函数空间拓扑}
\label{b1:subsec:200104}

先在 \(\mathcal D_K\) 上使用半范数 \(p_{K,m}\) 所给的拓扑。其收敛就是各阶导数一致收敛，可由度量
\[
 \mathrm{d}_K(\varphi,\psi)=\sum_{m=0}^\infty2^{-m-1}\min\{1,p_{K,m}(\varphi-\psi)\}
\]
描述。这个空间完备。事实上，Cauchy 列的每个导数都有一致极限 \(g_\alpha\)。在包含于 \(\Omega\) 的小球中，将普通积分恒等式
\[
 D^\alpha\varphi_j(x+he_i)-D^\alpha\varphi_j(x)
 =\int_0^h D^{\alpha+e_i}\varphi_j(x+te_i)\,\mathrm dt
\]
传到极限，得到 \(\partial_i g_\alpha=g_{\alpha+e_i}\)。故 \(g_0\) 光滑，其支集仍在 \(K\)，且原列在每个半范数中收敛到它。

\ProseParagraphBreak

整个 \(\mathcal D(\Omega)\) 是这些固定支集空间的并，却不能直接赋予 \(C^\infty(\Omega)\) 的局部一致拓扑。例如，在 \(\Omega=\mathbb R^d\) 中，一个固定小凸起向无穷远平移，在每个紧集上终将为零，但其积分始终不变。若这样的列在测试空间中趋于零，积分泛函就会失去连续性。

\ProseParagraphBreak

以下从半范数族构造完整局部凸拓扑的部分属于技术性补充；第一次阅读可以先接受随后\cref{b1:prop:test-sequence-convergence}与\cref{b1:prop:test-functional-continuity}给出的序列判据和有限阶估计，再回读这一构造。为避免把一个收敛列的判据误当成对全部开集的定义，设 \(\mathscr P\) 为 \(\mathcal D(\Omega)\) 上所有这样的半范数 \(q\)：对每个 \(K\Subset\Omega\)，存在 \(C>0\) 与整数 \(m\geq0\)，使
\[
 q(\varphi)\leq C\cdot p_{K,m}(\varphi)\quad(\varphi\in\mathcal D_K).
\]
以有限个 \(q\in\mathscr P\) 的小于正数的集合之交作为零点邻域基，定义 \(\mathcal D(\Omega)\) 的拓扑。半范数的不等式保证加法、数乘连续，而且每个 \(\mathcal D_K\) 的包含映射连续。由于这些邻域都是凸的，这是一种局部凸拓扑；它也是使各包含映射连续的最强局部凸拓扑。这最后一句只须注意：任何另一个由半范数给出的此类拓扑，其定义半范数在每个 \(\mathcal D_K\) 上连续，因而都已收入 \(\mathscr P\)。

\begin{proposition}\label{b1:prop:test-sequence-convergence}
序列 \(\varphi_j\) 在 \(\mathcal D(\Omega)\) 中趋于零，当且仅当存在同一个 \(K\Subset\Omega\) 包含所有 \(\operatorname{supp}\varphi_j\)，并且对每个 \(m\geq0\)，有 \(p_{K,m}(\varphi_j)\to0\)。
\end{proposition}
\begin{proof}
若具有右侧性质，则每个 \(q\in\mathscr P\) 的固定紧集估计给出 \(q(\varphi_j)\to0\)，所以在所定义拓扑中收敛。

反之，假定 \(\varphi_j\to0\)。全空间上各阶导数的上确界半范数属于 \(\mathscr P\)，故这些量趋于零。还须证明支集可统一限制。若不能，对上面紧集穷竭可选一个子列 \(\varphi_{j_l}\) 及点 \(x_l\notin K_l\)，使 \(\varphi_{j_l}(x_l)\ne0\)。每个有限的前缀有共同紧支集，因此指标可递增选取。点族 \((x_l)\) 在 \(\Omega\) 中局部有限。于是
\[
 q(\varphi)=\sum_l\frac{|\varphi(x_l)|}{|\varphi_{j_l}(x_l)|}
\]
是良定义的半范数：每个紧支集只遇到有限个 \(x_l\)，在固定 \(\mathcal D_K\) 上有 \(q\leq C_K\cdot p_{K,0}\)。故 \(q\in\mathscr P\)，但 \(q(\varphi_{j_l})\geq1\)，与收敛矛盾。这证明共同紧支集存在。
\end{proof}

\begin{proposition}\label{b1:prop:test-functional-continuity}
设 \(T:\mathcal D(\Omega)\rightarrow\mathbb K\) 线性。下列三个条件等价：
\begin{equivalences}
\item\label{b1:item:test-functional-topological}泛函 \(T\) 在上述拓扑中连续；
\item\label{b1:item:test-functional-sequential}每当 \(\varphi_j\to0\) 于 \(\mathcal D(\Omega)\)，都有 \(T\varphi_j\to0\)；
\item\label{b1:item:test-functional-order}对每个 \(K\Subset\Omega\)，存在 \(C_K>0\)、整数 \(m_K\geq0\)，使
\[
 |T\varphi|\leq C_K\cdot p_{K,m_K}(\varphi)\quad(\varphi\in\mathcal D_K).
\]
\end{equivalences}
\end{proposition}
\begin{proof}
连续性蕴涵序列连续性。若第三项对某个 \(K\) 失败，则对每个 \(j\geq1\) 可取 \(u_j\in\mathcal D_K\)，满足 \(|Tu_j|>j\cdot p_{K,j}(u_j)\)。令 \(v_j=u_j/|Tu_j|\)，则对每个固定的 \(m\)，在 \(j\geq m\) 时有 \(p_{K,m}(v_j)<1/j\)，但 \(|Tv_j|=1\)。这与第二项矛盾。最后，第三项说明半范数 \(\varphi\mapsto|T\varphi|\) 属于 \(\mathscr P\)，从而 \(T\) 连续。
\end{proof}

这个有限阶估计是后面使用测试函数拓扑的主要接口。微分和乘以固定光滑函数都保持测试列收敛；在线性泛函上，这足以证明转移过去的相应操作仍连续。紧集改变时允许估计中的阶数和常数一起改变，没有理由要求一个无界开集上的所有位置受同一有限阶控制。

% !TeX root = ../../Book1.tex
\section{分布}
\label{b1:sec:2002}

一个函数可以通过与测试函数的积分来辨认；分部积分又能把求导转移到测试函数上。沿着这个方向，研究对象不必先在每一点有函数值，只需规定它怎样作用于测试函数。这样得到的空间既包含局部可积函数，也容纳点质量及其任意阶导数。

% 实读 Hunter--Nachtergaele 作者官方 ch11.pdf，印刷291--292、294--297页：
% 11.2 的对偶配对、例11.5--11.6，11.3 的光滑乘法、定义11.10与例11.14。
% 原书这一段主要讨论 Schwartz 测试空间及温和分布；294页明确另有 D'。
% 本节改用任意开集上的 D，并自足展开局部有限阶、限制、支集和局部唯一性；
% 不移用原书的无穷远多项式增长条件，不依赖尚未建立的20.3光滑化。
% 中文异名据 Schwartz2010DistributionsChinese 高等教育出版社目录第一章核定；
% 本轮未读该中译本正文，目录只作为名称及伴读位置的证据，不承担数学证明。

\subsection{分布的定义}
\label{b1:subsec:200201}

本节固定开集 \(\Omega\subseteq\mathbb R^d\)，默认测试函数取复值，沿用 \(\mathcal D(\Omega)=C_c^\infty(\Omega)\)。光滑函数的偏导数记作 \(\partial^\alpha\)，与上一节的 \(D^\alpha\) 相同。

\begin{definition}\label{b1:def:distribution}
设 \(T:\mathcal D(\Omega)\rightarrow\mathbb C\) 为复线性泛函。若每当 \(\varphi_j\rightarrow0\) 为测试函数收敛列时都有 \(T(\varphi_j)\rightarrow0\)，则称 \(T\) 为 \(\Omega\) 上的\term[fenbu]{分布}[distribution]。所有这样的泛函组成向量空间 \(\mathcal D'(\Omega)\)，配对记作
\[
 \langle T,\varphi\rangle=T(\varphi).
\]
配对对两个变量均线性，不对测试函数取复共轭。
\end{definition}
\symidx{\(\mathcal D'(\Omega)\)：分布空间}

分布也称广义函数；这一中文名称见 Schwartz 的《广义函数论》\cite[第一章]{Schwartz2010DistributionsChinese}\termidx[guangyi-hanshu]{广义函数}\foreignidx{generalized function}。这里的“广义”不是给每个点补一个广义数值，而是改变描述对象的方式。实值理论采用实测试函数和实线性泛函，以下证明不变。

由\cref{b1:prop:test-functional-continuity}，上述连续性等价于：对每个紧集 \(K\subseteq\Omega\)，存在常数 \(C_K\) 和非负整数 \(m_K\)，使
\begin{equation}\label{b1:eq:distribution-local-order}
 |\langle T,\varphi\rangle|
 \leq C_K p_{K,m_K}(\varphi)
 \quad(\varphi\in\mathcal D_K).
\end{equation}
这就是分布的局部有限阶估计。阶数允许随 \(K\) 改变；若能对所有紧集选用同一个整数 \(m\)，则称这个\term[fenbu-de-jie]{分布的阶}[order of a distribution]不超过 \(m\)。局部估计并不要求 \(\Omega\) 上存在一个统一的阶数或常数。

\begin{definition}\label{b1:def:distribution-convergence}
设 \(T_j,T\in\mathcal D'(\Omega)\)。若对每个固定的 \(\varphi\in\mathcal D(\Omega)\) 都有
\[
 \langle T_j,\varphi\rangle\longrightarrow\langle T,\varphi\rangle,
\]
则称 \(T_j\) 在分布意义下收敛于 \(T\)，记作 \(T_j\rightarrow T\) 于 \(\mathcal D'(\Omega)\)。
\end{definition}

这里变化的是分布而不是测试函数；它与定义分布时要求的测试函数列连续性是两个不同的量词安排。

\begin{definition}\label{b1:def:distribution-support}
对开集 \(U\subseteq\Omega\)，将 \(\varphi\in\mathcal D(U)\) 在 \(\Omega\setminus U\) 上延为零，所得测试函数记为 \(\widetilde\varphi\)。规定
\[
 \langle T|_U,\varphi\rangle=\langle T,\widetilde\varphi\rangle,
\]
称 \(T|_U\) 为 \(T\) 在 \(U\) 上的限制。若 \(T|_U=0\)，则称 \(T\) 在 \(U\) 上为零。定义\term[fenbu-zhiji]{分布的支集}[support of a distribution]为
\[
 \operatorname{supp}T
 =\Omega\setminus\bigcup\{\,U\subseteq\Omega:U\text{ open},\ T|_U=0\,\}.
\]
\end{definition}
\symidx{\(\operatorname{supp}T\)：分布的支集}

\begin{proposition}\label{b1:prop:distribution-locality}
限制 \(T|_U\) 是 \(U\) 上的分布，并与进一步限制相容。支集是 \(\Omega\) 中的相对闭集；若测试函数的支集避开 \(\operatorname{supp}T\)，则 \(T\) 在该测试函数上的值为零。特别地，两个分布在某个开覆盖的每个成员上相等，就在 \(\Omega\) 上相等。
\end{proposition}
\begin{proof}
紧支集测试函数在 \(U\) 的边界附近已经为零，因而延零后仍光滑，且延零保持共同紧支集及所有阶导数的一致收敛。这证明限制的连续性；复线性和限制的相容性直接由定义得到。

记定义中零开集的并为 \(V\)。若 \(\varphi\in\mathcal D(V)\)，其紧支集被有限个这样的开集 \(U_1,\ldots,U_N\) 覆盖。由\cref{b1:thm:smooth-partition-unity}，可将它写成有限和 \(\varphi=\sum_{\ell=1}^N\varphi_\ell\)，其中 \(\varphi_\ell\in\mathcal D(U_\ell)\)。每项配对为零，所以 \(T|_V=0\)。支集外测试的结论随即成立。对两个分布之差作同一论证，就得到开覆盖上的唯一性。
\end{proof}

支集的意义还在于局部化：如果两个测试函数在 \(\operatorname{supp}T\) 的某个邻域相同，它们之差的支集便避开 \(\operatorname{supp}T\)，所以配对相同。对紧支集分布，这使我们能够扩大测试函数的范围。

\begin{corollary}\label{b1:cor:compact-distribution-smooth-action}
若 \(S=\operatorname{supp}T\) 是包含于 \(\Omega\) 的紧集，则 \(T\) 有有限阶。选取在 \(S\) 的邻域为 \(1\) 的 \(\chi\in\mathcal D(\Omega)\)，可以对任意 \(f\in C^\infty(\Omega)\) 规定
\[
 \langle T,f\rangle=\langle T,\chi f\rangle.
\]
这一规定不依赖 \(\chi\)，并与原来在测试函数上的作用一致。存在固定紧集 \(L\subseteq\Omega\)、整数 \(m\) 和常数 \(C\)，使
\[
 |\langle T,f\rangle|
 \leq C\max_{|\alpha|\leq m}\sup_{x\in L}|\partial^\alpha f(x)|.
\]
\end{corollary}
\begin{proof}
截断函数由\cref{b1:lem:smooth-cutoff}给出。若 \(\chi_1,\chi_2\) 都满足要求，则 \((\chi_1-\chi_2)f\) 为紧支集光滑函数，并在 \(S\) 的邻域为零，由上一命题，其配对为零。对原本紧支集的 \(f\)，同样将 \((1-\chi)f\) 配对为零，得到一致性。

在 \(L=\operatorname{supp}\chi\) 上应用\eqref{b1:eq:distribution-local-order}，再展开 \(\chi f\) 的有限阶导数，就得到所述估计。对任意测试函数 \(\varphi\)，\(L\) 上的导数上确界不超过 \(\operatorname{supp}\varphi\) 上相同阶导数的上确界，因为所有这些导数在该支集之外为零。因此这里的同一个 \(m\) 控制所有测试函数，证明 \(T\) 有有限阶。
\end{proof}

\subsection{正则分布}
\label{b1:subsec:200202}

\begin{definition}\label{b1:def:regular-distribution}
给定 \(f\in L^1_{\mathrm{loc}}(\Omega)\)，规定
\[
 \langle T_f,\varphi\rangle
 =\int_\Omega f(x)\varphi(x)\,\mathrm dx.
\]
由这种局部可积函数产生的分布称为\term[zhengze-fenbu]{正则分布}[regular distribution]。
\end{definition}

这个规定确实给出分布。若 \(\operatorname{supp}\varphi\subseteq K\)，则
\[
 |\langle T_f,\varphi\rangle|
 \leq\left(\int_K|f|\,\mathrm dx\right)p_{K,0}(\varphi).
\]
因此 \(T_f\) 为零阶分布。积分只看到几乎处处等价类；下面证明它也没有丢失等价类以外的信息。

\begin{proposition}\label{b1:prop:regular-distribution-injective}
映射 \(f\mapsto T_f\) 是从 \(L^1_{\mathrm{loc}}(\Omega)\) 到 \(\mathcal D'(\Omega)\) 的线性单射。若 \(f_j\rightarrow f\) 于 \(L^1_{\mathrm{loc}}\)，则 \(T_{f_j}\rightarrow T_f\) 于分布意义。
\end{proposition}
\begin{proof}
只需证明 \(T_f=0\) 蕴含 \(f=0\) 几乎处处。选取非负的 \(\rho\in C_c^\infty\bigl(B(0,1)\bigr)\)，使 \(\int\rho=1\)。这样的函数由上一节的光滑截断归一化得到。对 \(f\) 的一个 Lebesgue 点 \(x\in\Omega\)，当 \(r>0\) 足够小时，
\[
 \varphi_{x,r}(y)=r^{-d}\rho\left(\frac{y-x}r\right)
\]
属于 \(\mathcal D(\Omega)\)，所以 \(\int f\varphi_{x,r}=0\)。利用其积分为 \(1\)，有
\[
 \begin{aligned}
 |f(x)|
 &=\left|\int_\Omega\bigl(f(x)-f(y)\bigr)\varphi_{x,r}(y)\,\mathrm dy\right|\\
 &\leq\frac{\|\rho\|_\infty}{r^d}
        \int_{B(x,r)}|f(y)-f(x)|\,\mathrm dy
 \longrightarrow0.
 \end{aligned}
\]
由\cref{b1:thm:lebesgue-point}，几乎每一点都是这样的点，故单射性成立。最后，对固定测试函数及 \(K=\operatorname{supp}\varphi\)，
\[
 |\langle T_{f_j}-T_f,\varphi\rangle|
 \leq\|\varphi\|_\infty\cdot\|f_j-f\|_{L^1(K)}
 \longrightarrow0.
\]
\end{proof}

因此以后可以把 \(f\) 与 \(T_f\) 视为同一对象，但不能据此为一般分布写下点值 \(T(x)\)。同一命题在每个开子集上适用，故 \(T_f\) 在一个开集上为零，当且仅当 \(f\) 在那里几乎处处为零。这也说明分布支集对应的是函数的本质支集，而不是任选代表的非零点集合之闭包。

\subsection{Dirac 分布}
\label{b1:subsec:200203}

\begin{definition}\label{b1:def:dirac-distribution}
给定 \(a\in\Omega\)，以
\[
 \langle\delta_a,\varphi\rangle=\varphi(a)
\]
定义 \(a\) 处的\term[dirac-fenbu]{Dirac 分布}[Dirac distribution]。
\end{definition}

估计 \(|\varphi(a)|\leq p_{K,0}(\varphi)\) 对 \(a\in K\) 成立；若 \(a\notin K\)，配对为零。因此 \(\delta_a\) 是零阶分布。它与点质量测度作用于测试函数所得的泛函一致，但不是局部可积函数。

\begin{proposition}\label{b1:prop:dirac-not-regular}
Dirac 分布不是正则分布。更一般地，对任意多重指标 \(\alpha\)，泛函
\[
 \langle S_{a,\alpha},\varphi\rangle
 =(-1)^{|\alpha|}\partial^\alpha\varphi(a)
\]
是非正则分布，且其支集恰为 \(\{a\}\)。
\end{proposition}
\begin{proof}
点上导数由 \(p_{K,|\alpha|}\) 控制，所以 \(S_{a,\alpha}\) 连续。选取在原点附近为 \(1\) 的光滑截断 \(\chi\)，支集包含于 \(B(0,1)\)，令
\[
 \psi(z)=\frac{z^\alpha}{\alpha!}\chi(z),\quad
 \psi_r(x)=r^{|\alpha|}\psi\left(\frac{x-a}r\right).
\]
当 \(r\) 足够小时，\(\psi_r\in\mathcal D(\Omega)\)，且 \(\partial^\alpha\psi_r(a)=1\)。若 \(S_{a,\alpha}=T_g\)，则一方面配对的绝对值恒为 \(1\)，另一方面
\[
 |\langle T_g,\psi_r\rangle|
 \leq r^{|\alpha|}\|\psi\|_\infty
            \int_{B(a,r)}|g(x)|\,\mathrm dx
 \longrightarrow0,
\]
其中使用 \(g\) 的局部可积性及积分的绝对连续性，产生矛盾。取 \(\alpha=0\) 即包含 Dirac 分布的情形。

支集避开 \(a\) 的测试函数在 \(a\) 的邻域恒为零，其所有点上导数也为零；反过来，任意包含 \(a\) 的开集都容纳上述某个 \(\psi_r\)，其配对非零。由支集定义得到结论。
\end{proof}

Dirac 分布揭示了“零阶”与“正则”之间的区别：只需要函数值估计的泛函仍可能集中在一个 Lebesgue 零测集上。Hunter 与 Nachtergaele 的《Applied Analysis》\cite[Sections 11.2--11.3, pp. 291--296]{Hunter2001Applied}给出了点质量、正则分布和导数的同类构造；该书此处主要使用温和分布，本章的紧支集测试则不要求无穷远处的增长限制。

\subsection{分布导数}
\label{b1:subsec:200204}

若 \(f\) 足够光滑，反复分部积分会把它的偏导数转移到测试函数上，而紧支集消去边界项。这提示我们把右侧作为一般分布的求导规则。

\begin{definition}\label{b1:def:distributional-derivative}
设 \(T\in\mathcal D'(\Omega)\)，\(\alpha\) 为多重指标。定义其\term[fenbu-daoshu]{分布导数}[distributional derivative]为
\begin{equation}\label{b1:eq:distributional-derivative}
 \langle\partial^\alpha T,\varphi\rangle
 =(-1)^{|\alpha|}\langle T,\partial^\alpha\varphi\rangle.
\end{equation}
对 \(a\in C^\infty(\Omega)\)，定义光滑函数与分布的乘法为
\[
 \langle aT,\varphi\rangle=\langle T,a\varphi\rangle.
\]
\end{definition}
\symidx{\(\partial^\alpha T\)：分布导数}

\begin{proposition}\label{b1:prop:distribution-derivative-calculus}
上述导数与乘积都是分布。分布偏导数可交换，并满足
\[
 \partial^\alpha\partial^\beta T=\partial^{\alpha+\beta}T.
\]
求导与乘以固定光滑函数都保持分布收敛。若 \(f\in C^{|\alpha|}(\Omega)\)，则
\[
 \partial^\alpha T_f=T_{\partial^\alpha f}.
\]
此外，限制与求导、光滑乘法相容，且
\[
 \operatorname{supp}(\partial^\alpha T)\subseteq\operatorname{supp}T,\quad
 \operatorname{supp}(aT)\subseteq\operatorname{supp}a\cap\operatorname{supp}T.
\]
\end{proposition}
\begin{proof}
固定 \(K\)，若 \(T\) 满足阶数为 \(m\) 的局部估计，则
\[
 |\langle\partial^\alpha T,\varphi\rangle|
 \leq C_K p_{K,m+|\alpha|}(\varphi).
\]
光滑乘积的通常求导公式给出
\[
 p_{K,m}(a\varphi)
 \leq C_{K,m,a}p_{K,m}(\varphi),
\]
因为其中只出现 \(a\) 在 \(K\) 上有限多个有界导数。代入 \(T\) 的估计，就证明了 \(aT\) 的连续性。

对同一测试函数连续使用定义，得到
\[
 \langle\partial^\alpha\partial^\beta T,\varphi\rangle
 =(-1)^{|\alpha|+|\beta|}
       \langle T,\partial^{\alpha+\beta}\varphi\rangle,
\]
所以混合偏导数可交换。若 \(T_j\rightarrow T\)，分别用固定测试函数 \(\partial^\alpha\varphi\) 和 \(a\varphi\) 测试，即得对应导数与乘积的收敛。

对经典可微的 \(f\)，先将测试函数用光滑划分分为支集包含在 \(\Omega\) 内小矩形中的有限项。对每项使用 Fubini 定理和一维分部积分，反复 \(|\alpha|\) 次可得
\[
 \int_\Omega(\partial^\alpha f)\varphi\,\mathrm dx
 =(-1)^{|\alpha|}\int_\Omega f\,\partial^\alpha\varphi\,\mathrm dx.
\]
测试函数在每个矩形边界附近为零，故没有边界项。求和证明与经典导数一致。

延零与测试函数的求导、乘法相容，所以限制也相容。若 \(T\) 在某开集上为零，则其导数和 \(aT\) 在该开集上为零；若 \(a\) 在开集上恒为零，则 \(aT\) 也为零。支集的包含关系由此得到。
\end{proof}

\begin{proposition}[Leibniz 公式]\label{b1:prop:distribution-leibniz}
对 \(a\in C^\infty(\Omega)\) 与任意多重指标 \(\alpha\)，有
\[
 \partial^\alpha(aT)
 =\sum_{\beta\leq\alpha}\binom{\alpha}{\beta}
      (\partial^\beta a)\,\partial^{\alpha-\beta}T.
\]
\end{proposition}
\begin{proof}
先求一阶导数。对任意测试函数，
\[
 \begin{aligned}
 \langle\partial_i(aT),\varphi\rangle
 &=-\langle T,a\,\partial_i\varphi\rangle\\
 &=-\langle T,\partial_i(a\varphi)\rangle
       +\langle T,(\partial_i a)\varphi\rangle\\
 &=\langle a\,\partial_iT+(\partial_i a)T,\varphi\rangle.
 \end{aligned}
\]
再对 \(|\alpha|\) 归纳。给归纳式施加 \(\partial_i\)，每一项分别对光滑系数和分布求导；同一项的两个系数按
\(\binom{\alpha}{\beta}+\binom{\alpha}{\beta-e_i}
=\binom{\alpha+e_i}{\beta}\)
合并，越界系数规定为零，便得到下一阶公式。
\end{proof}

由定义，上一小节的点导数泛函正是 \(\partial^\alpha\delta_a\)，所以每个分布都能任意阶求导，但导数不必是正则分布。例如在一维有 \(x\delta_0=0\)，一阶 Leibniz 公式给出
\[
 x\delta_0'=-\delta_0.
\]
这些运算不依赖点值，却仍能给出精确的代数恒等式。对于一般两个分布，上述定义没有赋予它们乘法，因而不能把 \(\delta_0^2\) 等符号作为已经存在的对象使用。

局部有限阶也不等于全局有限阶：把越来越高阶的点导数放在趋于无穷远的位置，可以使每个紧集只遇到有限多项，却找不到统一的阶数界。具体构造留在习题\ref{b1:ex:20-locally-finite-unbounded-order}。

分布空间容纳这种随位置变化的阶数，同时又允许统一求导。本章的下一步是研究如何用光滑函数逼近分布；再把分布导数要求为局部可积函数，就导向\secref{b1:sec:2004}的弱导数。

% !TeX root = ../../Book1.tex
\section{卷积与光滑化}
\label{b1:sec:2003}

\subsecref{b1:subsec:130403}已经用可积核构造函数的逼近恒等。本节把平均推广到分布：把核的平移看成测试函数，由分布作用于它。所得对象虽然是普通光滑函数，其导数却能保留原分布的微分信息。整个构造先在远离边界的区域进行，不预设分布可以延拓到开集之外。

% 实读 Hunter2001Applied 作者官方 ch11.pdf，11.3--11.4，印刷 pp.297--301。
% 原书讨论 Schwartz 测试函数与缓增分布，11.23 只有稠密性证明概要；
% 本节一般开集 D' 的有限阶差商、共同紧支集 Riemann 和及局部化证明另行完整组织。
% 名称依据：Oleinik2008PDEChinese官方目录1.2用“磨光函数”；只核名称，未读该节定义。
% “磨光子”另见北京大学公开《实变函数》前言，来源记录限定到名称，不据此声称核对其归一化。

\subsection{磨光函数}
\label{b1:subsec:200301}

取非负偶函数 \(\rho\in C_c^\infty(\mathbb R^d)\)，使
\[
 \operatorname{supp}\rho\subset\overline{B(0,1)},
 \quad \int_{\mathbb R^d}\rho=1,
 \quad \rho_\varepsilon(x)\coloneqq\varepsilon^{-d}\cdot \rho(x/\varepsilon).
\]
满足上述条件的核称为本章的\term[moguang-hanshu]{磨光函数}[mollifier]。名称可对照 Oleinik 的《偏微分方程讲义》\cite[第1.2节]{Oleinik2008PDEChinese}；中文资料也使用“磨光子”这一名称。\termidx[moguangzi]{磨光子}\footnote{北京大学数学学院公开的《实变函数》前言，在介绍\chapref{b1:ch:04}的 \(L^p\) 与 Sobolev 空间工具时采用“磨光子”。该用词说明见\href{https://math101.pku.edu.cn/docs/2025-04/b74312fec0b54f0c90ff66270f2bded8.pdf}{学校公布的前言}；这里仅提示名称，不据前言推断其核函数的全部约定。}

\ProseParagraphBreak

例如可将 \(|x|<1\) 上的 \(\exp\bigl(-1/(1-|x|^2)\bigr)\) 在球外补零，再除以其正积分；光滑性已在\secref{b1:sec:2001}的截断构造中说明。这是\cref{b1:def:approximate-identity}中的一种特殊核，光滑性和紧支集是本节新增的要求，偶性则用来简化反射记号。

\ProseParagraphBreak

对开集 \(\Omega\subset\mathbb R^d\)，记
\[
 \Omega_\varepsilon
 \coloneqq\{\,x\in\Omega\mid\operatorname{dist}(x,\Omega^c)>\varepsilon\,\}.
\]
约定到空集的距离为无穷，故 \(\Omega=\mathbb R^d\) 时 \(\Omega_\varepsilon=\mathbb R^d\)。若 \(x\in\Omega_\varepsilon\)，则 \(y\mapsto\rho_\varepsilon(x-y)\) 的支集紧含于 \(\Omega\)；这才保证一般分布可以作用于它。

\subsection{分布卷积}
\label{b1:subsec:200302}

\begin{definition}\label{b1:def:distribution-test-convolution}
设 \(T\in\mathcal D'(\Omega)\)、\(\eta\in C_c^\infty(\mathbb R^d)\)，且 \(\operatorname{supp}\eta\subset\overline{B(0,a)}\)，其中 \(a>0\)。在 \(\Omega_a\) 上定义 \(T\) 与测试函数 \(\eta\) 的\term[fenbu-juanji]{分布卷积}[convolution of a distribution with a test function]为
\[
 (T*\eta)(x)\coloneqq\langle T_y,\eta(x-y)\rangle.
\]
下标 \(y\) 指明分布所作用的变量，不表示对 \(T\) 作逐点取值。
\end{definition}
\symidx{\(T*\eta\)：分布与测试函数的卷积}

若 \(T=T_f\) 来自 \(f\in L^1_{\mathrm{loc}}(\Omega)\)，该定义成为
\[
 (T_f*\eta)(x)=\int_\Omega f(y)\cdot\eta(x-y)\,\mathrm dy.
\]
被积函数支集紧含于 \(\Omega\)，所以每个 \(x\in\Omega_a\) 处都绝对可积，并与前文函数卷积一致。这里没有定义两个任意分布的卷积；即使两个常函数 \(1\) 的通常卷积，在全空间也会发散。紧支集等附加条件可以支持更广的分布卷积，但不参与本节证明。

\begin{proposition}\label{b1:prop:distribution-convolution-smooth}
上述 \(T*\eta\) 属于 \(C^\infty(\Omega_a)\)。对任意多重指标 \(\alpha\)，在 \(\Omega_a\) 上有
\[
 D^\alpha(T*\eta)=T*(D^\alpha\eta)=(D^\alpha T)*\eta.
\]
\end{proposition}
\begin{proof}
固定 \(K\Subset\Omega_a\)，取 \(K\) 的稍大紧邻域 \(K'\Subset\Omega_a\)。当 \(x\in K'\) 时，所有 \(\eta(x-\cdot)\) 的支集都在紧集 \(H=K'+\overline{B(0,a)}\Subset\Omega\) 中。由\eqref{b1:eq:distribution-local-order}，存在 \(C,m\) 使
\[
 |\langle T,\psi\rangle|\le C\cdot p_{H,m}(\psi)
 \quad(\psi\in\mathcal D_H).
\]
对每个坐标方向 \(e_j\)，Taylor 公式及 \(\eta\) 各阶导数的有界性给出
\[
 p_{H,m}\left(
 \frac{\eta(x+he_j-\cdot)-\eta(x-\cdot)}h
 -(\partial_j\eta)(x-\cdot)\right)\le C_\eta|h|
\]
对 \(x\in K\) 和充分小的 \(h\) 一致成立。将 \(T\) 作用于差商，便证明第一阶导数公式；同一估计也证明所得导数连续。把 \(\eta\) 换成其各阶导数，迭代即得全部光滑性与第一项等式。

分布导数的定义又给出
\[
 \langle D^\alpha T_y,\eta(x-y)\rangle
 =(-1)^{|\alpha|}\langle T_y,D_y^\alpha\eta(x-y)\rangle
 =\langle T_y,(D^\alpha\eta)(x-y)\rangle,
\]
因为对 \(y\) 微分还产生一个 \((-1)^{|\alpha|}\)。这证明第二项等式。
\end{proof}

例如在全空间，\(\delta_b*\eta=\eta(\cdot-b)\)，而 \((D^\alpha\delta_b)*\eta=D^\alpha\eta(\cdot-b)\)。点质量及其导数被变成普通光滑函数，但导数阶数没有丢失。

\subsection{正则化}
\label{b1:subsec:200303}

为了讨论收敛，还须把光滑函数 \(T*\eta\) 再与测试函数配对。记反射核 \(\check\eta(z)=\eta(-z)\)。反射不可由卷积记号中擅自省去；当 \(\eta=\rho_\varepsilon\) 时，偶性才使它等于原核。

\begin{proposition}\label{b1:prop:distribution-convolution-pairing}
若 \(\varphi\in\mathcal D(\Omega_a)\)，将其在外部补零，则
\[
 \langle T*\eta,\varphi\rangle
 =\langle T,\check\eta*\varphi\rangle.
\]
右端的卷积属于 \(\mathcal D(\Omega)\)。
\end{proposition}
\begin{proof}
令 \(K=\operatorname{supp}\varphi\)，则
\[
 q(y)=\int_{\mathbb R^d}\varphi(x)\cdot\eta(x-y)\,\mathrm dx
 =(\check\eta*\varphi)(y)
\]
光滑，且支集包含在 \(H=K+\overline{B(0,a)}\Subset\Omega\) 中。以下证明 \(T\) 与这个积分确能交换。

取包含 \(K\) 的有界闭矩形 \(Q\)，并将其分成网格趋于零的小矩形 \(Q_i\)，各选一点 \(x_i\)。函数值 Riemann 和
\[
 q_N(y)=\sum_i m_d(Q_i)\cdot\varphi(x_i)\cdot\eta(x_i-y)
\]
都属于 \(\mathcal D_H\)。对任意固定阶数 \(j\)，\(\varphi(x)D_y^\beta\eta(x-y)\) 在共同紧集上一致连续，\(|\beta|\le j\)，所以 Riemann 和误差关于 \(y\) 一致趋零，即 \(p_{H,j}(q_N-q)\to0\)。有限阶估计因此给出 \(\langle T,q_N\rangle\to\langle T,q\rangle\)。另一方面，由线性性，每个配对都是连续紧支集函数 \(x\mapsto\varphi(x)(T*\eta)(x)\) 的 Riemann 和；在 \(\Omega_a\) 外该乘积补零。它们趋于 \(\int\varphi(x)(T*\eta)(x)\,\mathrm dx\)，等式得证。
\end{proof}

\begin{theorem}[分布的局部光滑化]\label{b1:thm:distribution-regularization}
设 \(T\in\mathcal D'(\Omega)\)，令 \(T_\varepsilon=T*\rho_\varepsilon\in C^\infty(\Omega_\varepsilon)\)。则对每个固定的 \(\varphi\in\mathcal D(\Omega)\)，当 \(\varepsilon\) 充分小时配对有定义，且
\[
 \langle T_\varepsilon,\varphi\rangle\longrightarrow\langle T,\varphi\rangle.
\]
每个 \(D^\alpha T_\varepsilon\) 也以同样的局部意义趋于 \(D^\alpha T\)。
\end{theorem}
\begin{proof}
设 \(K=\operatorname{supp}\varphi\)，取 \(a>0\) 使 \(H=K+\overline{B(0,a)}\Subset\Omega\)。当 \(0<\varepsilon<a\) 时，\(\check\rho_\varepsilon*\varphi\) 和 \(\varphi\) 的支集均在 \(H\) 中，而且
\[
 D^\beta(\check\rho_\varepsilon*\varphi-\varphi)(y)
 =\int\rho(z)\cdot\bigl(D^\beta\varphi(y+\varepsilon z)-D^\beta\varphi(y)\bigr)\,\mathrm dz.
\]
因 \(D^\beta\varphi\) 一致连续、\(\rho\) 只在单位球中有质量，每个阶数的上确界范数都趋于零。前一命题与 \(T\) 的连续性给出所求配对极限。再将 \(T\) 换成 \(D^\alpha T\)，并用卷积微分公式，即得导数结论。
\end{proof}

这里 \(T_\varepsilon\) 的定义域随 \(\varepsilon\) 改变；“局部”指每个固定测试函数最终都能使用，不是把 \(T_\varepsilon\) 未经说明地当成整个 \(\Omega\) 上的函数。也不要求 \(T_\varepsilon\) 本身有逐点极限：例如 \(\delta_0*\rho_\varepsilon=\rho_\varepsilon\)，其峰值通常随 \(\varepsilon^{-d}\) 增长。

\subsection{逼近恒等}
\label{b1:subsec:200304}

\begin{proposition}\label{b1:prop:local-smoothing-approximation}
设 \(K\Subset\Omega\)。若 \(f\in L^p_{\mathrm{loc}}(\Omega)\)、\(1\le p<\infty\)，则
\[
 \|f*\rho_\varepsilon-f\|_{L^p(K)}\longrightarrow0.
\]
若 \(f\in C^k(\Omega)\)，则对所有 \(|\alpha|\le k\)，\(D^\alpha(f*\rho_\varepsilon)\to D^\alpha f\) 在 \(K\) 上一致成立。
\end{proposition}
\begin{proof}
由\secref{b1:sec:2001}的光滑截断，取 \(\chi\in C_c^\infty(\Omega)\)，使它在 \(K\) 的一个邻域上等于 \(1\)。将 \(\chi f\) 补零到全空间；当 \(\varepsilon\) 充分小时，它与 \(f\) 的卷积在 \(K\) 上相同。有限 \(p\) 的结论因此直接归结为\cref{b1:prop:approximate-identity-lp}。

若 \(f\in C^k\)，则 \(\chi f\) 的各阶导数都是连续紧支集函数。对每个 \(|\alpha|\le k\)，卷积微分公式把所需差转为 \((D^\alpha(\chi f))*\rho_\varepsilon-D^\alpha(\chi f)\)，再用同一命题的一致收敛部分即可。
\end{proof}

不能把有限 \(p\) 的结论照搬到 \(p=\infty\)。例如 \(f=\mathbf1_{(0,\infty)}\) 在一维原点附近有跳跃，偶核卷积在原点取值 \(1/2\)，并在附近连续，故其与 \(f\) 的本质上确界距离至少为 \(1/2\)。局部分布收敛、有限 \(p\) 范数收敛和光滑函数的逐阶一致收敛是不同层次，后续使用时须说明所需的是哪一种。

卷积平均怎样消除跳跃并保持局部信息，可见\cref{b1:fig:visual-f15}。
% F15：本书原创矢量示意；依据所在节的定义、证明或明确模型。
\begin{figure}[htbp]
\centering
\begin{tikzpicture}[x=1cm,y=1cm,>=stealth,every node/.style={font=\small},line width=.65pt]
\node[] at (6,3.4) {单位积分光滑核的卷积逼近};
\draw[->,gray] (0.88,0)--(11.8,0) node[right] {\(x\)};\draw[->,gray] (1,-0.12)--(1,2.7);
\draw[gray,dashed] plot coordinates {(3.0000,0.0000) (3.0000,2.2000) (7.0000,2.2000) (7.0000,0.0000)};
\draw[AnalysisBlue,thick] plot coordinates {(0.8000,0.0000) (0.8384,0.0000) (0.8767,0.0000) (0.9151,0.0000) (0.9534,0.0000) (0.9918,0.0000) (1.0301,0.0000) (1.0685,0.0000) (1.1068,0.0000) (1.1452,0.0000) (1.1836,0.0000) (1.2219,0.0000) (1.2603,0.0000) (1.2986,0.0000) (1.3370,0.0000) (1.3753,0.0000) (1.4137,0.0000) (1.4521,0.0000) (1.4904,0.0000) (1.5288,0.0001) (1.5671,0.0005) (1.6055,0.0018) (1.6438,0.0044) (1.6822,0.0087) (1.7205,0.0150) (1.7589,0.0236) (1.7973,0.0344) (1.8356,0.0474) (1.8740,0.0627) (1.9123,0.0801) (1.9507,0.0997) (1.9890,0.1212) (2.0274,0.1446) (2.0658,0.1698) (2.1041,0.1967) (2.1425,0.2251) (2.1808,0.2550) (2.2192,0.2863) (2.2575,0.3189) (2.2959,0.3527) (2.3342,0.3877) (2.3726,0.4236) (2.4110,0.4605) (2.4493,0.4983) (2.4877,0.5369) (2.5260,0.5762) (2.5644,0.6163) (2.6027,0.6569) (2.6411,0.6981) (2.6795,0.7398) (2.7178,0.7819) (2.7562,0.8244) (2.7945,0.8672) (2.8329,0.9103) (2.8712,0.9536) (2.9096,0.9971) (2.9479,1.0407) (2.9863,1.0844) (3.0247,1.1281) (3.0630,1.1718) (3.1014,1.2153) (3.1397,1.2588) (3.1781,1.3020) (3.2164,1.3451) (3.2548,1.3878) (3.2932,1.4302) (3.3315,1.4722) (3.3699,1.5138) (3.4082,1.5548) (3.4466,1.5952) (3.4849,1.6351) (3.5233,1.6742) (3.5616,1.7126) (3.6000,1.7501) (3.6384,1.7868) (3.6767,1.8224) (3.7151,1.8570) (3.7534,1.8905) (3.7918,1.9227) (3.8301,1.9537) (3.8685,1.9832) (3.9068,2.0112) (3.9452,2.0376) (3.9836,2.0623) (4.0219,2.0852) (4.0603,2.1061) (4.0986,2.1251) (4.1370,2.1419) (4.1753,2.1566) (4.2137,2.1690) (4.2521,2.1791) (4.2904,2.1870) (4.3288,2.1927) (4.3671,2.1965) (4.4055,2.1987) (4.4438,2.1997) (4.4822,2.2000) (4.5205,2.2000) (4.5589,2.2000) (4.5973,2.2000) (4.6356,2.2000) (4.6740,2.2000) (4.7123,2.2000) (4.7507,2.2000) (4.7890,2.2000) (4.8274,2.2000) (4.8658,2.2000) (4.9041,2.2000) (4.9425,2.2000) (4.9808,2.2000) (5.0192,2.2000) (5.0575,2.2000) (5.0959,2.2000) (5.1342,2.2000) (5.1726,2.2000) (5.2110,2.2000) (5.2493,2.2000) (5.2877,2.2000) (5.3260,2.2000) (5.3644,2.2000) (5.4027,2.2000) (5.4411,2.2000) (5.4795,2.2000) (5.5178,2.2000) (5.5562,2.1997) (5.5945,2.1987) (5.6329,2.1965) (5.6712,2.1927) (5.7096,2.1870) (5.7479,2.1791) (5.7863,2.1690) (5.8247,2.1566) (5.8630,2.1419) (5.9014,2.1251) (5.9397,2.1061) (5.9781,2.0852) (6.0164,2.0623) (6.0548,2.0376) (6.0932,2.0112) (6.1315,1.9832) (6.1699,1.9537) (6.2082,1.9227) (6.2466,1.8905) (6.2849,1.8570) (6.3233,1.8224) (6.3616,1.7868) (6.4000,1.7501) (6.4384,1.7126) (6.4767,1.6742) (6.5151,1.6351) (6.5534,1.5952) (6.5918,1.5548) (6.6301,1.5138) (6.6685,1.4722) (6.7068,1.4302) (6.7452,1.3878) (6.7836,1.3451) (6.8219,1.3020) (6.8603,1.2588) (6.8986,1.2153) (6.9370,1.1718) (6.9753,1.1281) (7.0137,1.0844) (7.0521,1.0407) (7.0904,0.9971) (7.1288,0.9536) (7.1671,0.9103) (7.2055,0.8672) (7.2438,0.8244) (7.2822,0.7819) (7.3205,0.7398) (7.3589,0.6981) (7.3973,0.6569) (7.4356,0.6163) (7.4740,0.5762) (7.5123,0.5369) (7.5507,0.4983) (7.5890,0.4605) (7.6274,0.4236) (7.6658,0.3877) (7.7041,0.3527) (7.7425,0.3189) (7.7808,0.2863) (7.8192,0.2550) (7.8575,0.2251) (7.8959,0.1967) (7.9342,0.1698) (7.9726,0.1446) (8.0110,0.1212) (8.0493,0.0997) (8.0877,0.0801) (8.1260,0.0627) (8.1644,0.0474) (8.2027,0.0344) (8.2411,0.0236) (8.2795,0.0150) (8.3178,0.0087) (8.3562,0.0044) (8.3945,0.0018) (8.4329,0.0005) (8.4712,0.0001) (8.5096,0.0000) (8.5479,0.0000) (8.5863,-0.0000) (8.6247,0.0000) (8.6630,0.0000) (8.7014,0.0000) (8.7397,0.0000) (8.7781,0.0000) (8.8164,0.0000) (8.8548,0.0000) (8.8932,0.0000) (8.9315,0.0000) (8.9699,0.0000) (9.0082,0.0000) (9.0466,0.0000) (9.0849,0.0000) (9.1233,0.0000) (9.1616,0.0000) (9.2000,0.0000)};
\node[AnalysisBlue] at (10,2.4) {\(\varepsilon=0.8\)};
\draw[orange!80!black,thick] plot coordinates {(0.8000,0.0000) (0.8384,0.0000) (0.8767,0.0000) (0.9151,0.0000) (0.9534,0.0000) (0.9918,0.0000) (1.0301,0.0000) (1.0685,0.0000) (1.1068,0.0000) (1.1452,0.0000) (1.1836,0.0000) (1.2219,0.0000) (1.2603,0.0000) (1.2986,0.0000) (1.3370,0.0000) (1.3753,0.0000) (1.4137,0.0000) (1.4521,0.0000) (1.4904,0.0000) (1.5288,0.0000) (1.5671,0.0000) (1.6055,0.0000) (1.6438,0.0000) (1.6822,0.0000) (1.7205,0.0000) (1.7589,0.0000) (1.7973,0.0000) (1.8356,0.0000) (1.8740,0.0000) (1.9123,0.0000) (1.9507,0.0000) (1.9890,0.0000) (2.0274,0.0000) (2.0658,0.0000) (2.1041,0.0000) (2.1425,0.0000) (2.1808,0.0000) (2.2192,0.0000) (2.2575,0.0000) (2.2959,0.0012) (2.3342,0.0069) (2.3726,0.0203) (2.4110,0.0425) (2.4493,0.0737) (2.4877,0.1133) (2.5260,0.1606) (2.5644,0.2148) (2.6027,0.2750) (2.6411,0.3405) (2.6795,0.4106) (2.7178,0.4847) (2.7562,0.5621) (2.7945,0.6423) (2.8329,0.7248) (2.8712,0.8091) (2.9096,0.8949) (2.9479,0.9816) (2.9863,1.0688) (3.0247,1.1562) (3.0630,1.2433) (3.1014,1.3297) (3.1397,1.4151) (3.1781,1.4990) (3.2164,1.5809) (3.2548,1.6603) (3.2932,1.7368) (3.3315,1.8098) (3.3699,1.8787) (3.4082,1.9428) (3.4466,2.0014) (3.4849,2.0537) (3.5233,2.0989) (3.5616,2.1361) (3.6000,2.1648) (3.6384,2.1844) (3.6767,2.1954) (3.7151,2.1994) (3.7534,2.2000) (3.7918,2.2000) (3.8301,2.2000) (3.8685,2.2000) (3.9068,2.2000) (3.9452,2.2000) (3.9836,2.2000) (4.0219,2.2000) (4.0603,2.2000) (4.0986,2.2000) (4.1370,2.2000) (4.1753,2.2000) (4.2137,2.2000) (4.2521,2.2000) (4.2904,2.2000) (4.3288,2.2000) (4.3671,2.2000) (4.4055,2.2000) (4.4438,2.2000) (4.4822,2.2000) (4.5205,2.2000) (4.5589,2.2000) (4.5973,2.2000) (4.6356,2.2000) (4.6740,2.2000) (4.7123,2.2000) (4.7507,2.2000) (4.7890,2.2000) (4.8274,2.2000) (4.8658,2.2000) (4.9041,2.2000) (4.9425,2.2000) (4.9808,2.2000) (5.0192,2.2000) (5.0575,2.2000) (5.0959,2.2000) (5.1342,2.2000) (5.1726,2.2000) (5.2110,2.2000) (5.2493,2.2000) (5.2877,2.2000) (5.3260,2.2000) (5.3644,2.2000) (5.4027,2.2000) (5.4411,2.2000) (5.4795,2.2000) (5.5178,2.2000) (5.5562,2.2000) (5.5945,2.2000) (5.6329,2.2000) (5.6712,2.2000) (5.7096,2.2000) (5.7479,2.2000) (5.7863,2.2000) (5.8247,2.2000) (5.8630,2.2000) (5.9014,2.2000) (5.9397,2.2000) (5.9781,2.2000) (6.0164,2.2000) (6.0548,2.2000) (6.0932,2.2000) (6.1315,2.2000) (6.1699,2.2000) (6.2082,2.2000) (6.2466,2.2000) (6.2849,2.1994) (6.3233,2.1954) (6.3616,2.1844) (6.4000,2.1648) (6.4384,2.1361) (6.4767,2.0989) (6.5151,2.0537) (6.5534,2.0014) (6.5918,1.9428) (6.6301,1.8787) (6.6685,1.8098) (6.7068,1.7368) (6.7452,1.6603) (6.7836,1.5809) (6.8219,1.4990) (6.8603,1.4151) (6.8986,1.3297) (6.9370,1.2433) (6.9753,1.1562) (7.0137,1.0688) (7.0521,0.9816) (7.0904,0.8949) (7.1288,0.8091) (7.1671,0.7248) (7.2055,0.6423) (7.2438,0.5621) (7.2822,0.4847) (7.3205,0.4106) (7.3589,0.3405) (7.3973,0.2750) (7.4356,0.2148) (7.4740,0.1606) (7.5123,0.1133) (7.5507,0.0737) (7.5890,0.0425) (7.6274,0.0203) (7.6658,0.0069) (7.7041,0.0012) (7.7425,0.0000) (7.7808,0.0000) (7.8192,0.0000) (7.8575,0.0000) (7.8959,0.0000) (7.9342,0.0000) (7.9726,0.0000) (8.0110,0.0000) (8.0493,0.0000) (8.0877,0.0000) (8.1260,0.0000) (8.1644,0.0000) (8.2027,0.0000) (8.2411,0.0000) (8.2795,0.0000) (8.3178,0.0000) (8.3562,0.0000) (8.3945,0.0000) (8.4329,0.0000) (8.4712,0.0000) (8.5096,0.0000) (8.5479,0.0000) (8.5863,0.0000) (8.6247,0.0000) (8.6630,0.0000) (8.7014,0.0000) (8.7397,0.0000) (8.7781,0.0000) (8.8164,0.0000) (8.8548,0.0000) (8.8932,0.0000) (8.9315,0.0000) (8.9699,0.0000) (9.0082,0.0000) (9.0466,0.0000) (9.0849,0.0000) (9.1233,0.0000) (9.1616,0.0000) (9.2000,0.0000)};
\node[orange!80!black] at (10,1.7999999999999998) {\(\varepsilon=0.4\)};
\draw[black,thick] plot coordinates {(0.8000,0.0000) (0.8384,0.0000) (0.8767,0.0000) (0.9151,0.0000) (0.9534,0.0000) (0.9918,0.0000) (1.0301,0.0000) (1.0685,0.0000) (1.1068,0.0000) (1.1452,0.0000) (1.1836,0.0000) (1.2219,0.0000) (1.2603,0.0000) (1.2986,0.0000) (1.3370,0.0000) (1.3753,0.0000) (1.4137,0.0000) (1.4521,0.0000) (1.4904,0.0000) (1.5288,0.0000) (1.5671,0.0000) (1.6055,0.0000) (1.6438,0.0000) (1.6822,0.0000) (1.7205,0.0000) (1.7589,0.0000) (1.7973,0.0000) (1.8356,0.0000) (1.8740,0.0000) (1.9123,0.0000) (1.9507,0.0000) (1.9890,0.0000) (2.0274,0.0000) (2.0658,0.0000) (2.1041,0.0000) (2.1425,0.0000) (2.1808,0.0000) (2.2192,0.0000) (2.2575,0.0000) (2.2959,0.0000) (2.3342,0.0000) (2.3726,0.0000) (2.4110,0.0000) (2.4493,0.0000) (2.4877,0.0000) (2.5260,0.0000) (2.5644,0.0000) (2.6027,0.0000) (2.6411,0.0000) (2.6795,0.0000) (2.7178,0.0000) (2.7562,0.0113) (2.7945,0.0762) (2.8329,0.2000) (2.8712,0.3684) (2.9096,0.5677) (2.9479,0.7869) (2.9863,1.0168) (3.0247,1.2495) (3.0630,1.4772) (3.1014,1.6917) (3.1397,1.8834) (3.1781,2.0406) (3.2164,2.1488) (3.2548,2.1959) (3.2932,2.2000) (3.3315,2.2000) (3.3699,2.2000) (3.4082,2.2000) (3.4466,2.2000) (3.4849,2.2000) (3.5233,2.2000) (3.5616,2.2000) (3.6000,2.2000) (3.6384,2.2000) (3.6767,2.2000) (3.7151,2.2000) (3.7534,2.2000) (3.7918,2.2000) (3.8301,2.2000) (3.8685,2.2000) (3.9068,2.2000) (3.9452,2.2000) (3.9836,2.2000) (4.0219,2.2000) (4.0603,2.2000) (4.0986,2.2000) (4.1370,2.2000) (4.1753,2.2000) (4.2137,2.2000) (4.2521,2.2000) (4.2904,2.2000) (4.3288,2.2000) (4.3671,2.2000) (4.4055,2.2000) (4.4438,2.2000) (4.4822,2.2000) (4.5205,2.2000) (4.5589,2.2000) (4.5973,2.2000) (4.6356,2.2000) (4.6740,2.2000) (4.7123,2.2000) (4.7507,2.2000) (4.7890,2.2000) (4.8274,2.2000) (4.8658,2.2000) (4.9041,2.2000) (4.9425,2.2000) (4.9808,2.2000) (5.0192,2.2000) (5.0575,2.2000) (5.0959,2.2000) (5.1342,2.2000) (5.1726,2.2000) (5.2110,2.2000) (5.2493,2.2000) (5.2877,2.2000) (5.3260,2.2000) (5.3644,2.2000) (5.4027,2.2000) (5.4411,2.2000) (5.4795,2.2000) (5.5178,2.2000) (5.5562,2.2000) (5.5945,2.2000) (5.6329,2.2000) (5.6712,2.2000) (5.7096,2.2000) (5.7479,2.2000) (5.7863,2.2000) (5.8247,2.2000) (5.8630,2.2000) (5.9014,2.2000) (5.9397,2.2000) (5.9781,2.2000) (6.0164,2.2000) (6.0548,2.2000) (6.0932,2.2000) (6.1315,2.2000) (6.1699,2.2000) (6.2082,2.2000) (6.2466,2.2000) (6.2849,2.2000) (6.3233,2.2000) (6.3616,2.2000) (6.4000,2.2000) (6.4384,2.2000) (6.4767,2.2000) (6.5151,2.2000) (6.5534,2.2000) (6.5918,2.2000) (6.6301,2.2000) (6.6685,2.2000) (6.7068,2.2000) (6.7452,2.1959) (6.7836,2.1488) (6.8219,2.0406) (6.8603,1.8834) (6.8986,1.6917) (6.9370,1.4772) (6.9753,1.2495) (7.0137,1.0168) (7.0521,0.7869) (7.0904,0.5677) (7.1288,0.3684) (7.1671,0.2000) (7.2055,0.0762) (7.2438,0.0113) (7.2822,0.0000) (7.3205,0.0000) (7.3589,0.0000) (7.3973,0.0000) (7.4356,0.0000) (7.4740,0.0000) (7.5123,0.0000) (7.5507,0.0000) (7.5890,0.0000) (7.6274,0.0000) (7.6658,0.0000) (7.7041,0.0000) (7.7425,0.0000) (7.7808,0.0000) (7.8192,0.0000) (7.8575,0.0000) (7.8959,0.0000) (7.9342,0.0000) (7.9726,0.0000) (8.0110,0.0000) (8.0493,0.0000) (8.0877,0.0000) (8.1260,0.0000) (8.1644,0.0000) (8.2027,0.0000) (8.2411,0.0000) (8.2795,0.0000) (8.3178,0.0000) (8.3562,0.0000) (8.3945,0.0000) (8.4329,0.0000) (8.4712,0.0000) (8.5096,0.0000) (8.5479,0.0000) (8.5863,0.0000) (8.6247,0.0000) (8.6630,0.0000) (8.7014,0.0000) (8.7397,0.0000) (8.7781,0.0000) (8.8164,0.0000) (8.8548,0.0000) (8.8932,0.0000) (8.9315,0.0000) (8.9699,0.0000) (9.0082,0.0000) (9.0466,0.0000) (9.0849,0.0000) (9.1233,0.0000) (9.1616,0.0000) (9.2000,0.0000)};
\node[black] at (10,1.2) {\(\varepsilon=0.15\)};

\end{tikzpicture}
\caption[卷积如何移动、平均与光滑化]{卷积如何移动、平均与光滑化。虚线为区间示性函数；实线为它与尺度化光滑核的卷积。核取 \(C\exp(-1/(1-x^2))\) 于 \(|x|<1\)，区间外为零，其中常数使积分为一。}
\label{b1:fig:visual-f15}
\end{figure}


% !TeX root = ../../Book1.tex
\section{弱导数}
\label{b1:sec:2004}

分布总可以求导，但导数未必还能由函数表示。弱导数研究的正是后一个问题：给定局部可积函数，把它作为正则分布求导后，结果是否仍是一个局部可积函数？这个要求比几乎处处存在经典导数更强，因为分部积分还会检测跳跃以及集中在零测集上的变化。

% 实读：Hunter2001Applied，作者官方ch12.pdf，§12.9印刷362--364页，
% 定义12.63--12.68；本节主要对照12.64的多指标弱导数与L1loc范围。
% 原PDF文本层乱码，实际读作者公开原页图像，不声称这些页含本节全部证明。
% 唯一性承接2002正则嵌入；局部Lipschitz证明用16章Rademacher及差商平移；
% 一维代表用零积分测试函数原函数、Fubini及1604绝对连续微积分基本定理自足证明。

\subsection{定义}
\label{b1:subsec:200401}

设 \(\Omega\subseteq\mathbb R^d\) 为开集，函数可取实值或复值。沿用多指标记号 \(\alpha=(\alpha_1,\ldots,\alpha_d)\in\mathbb N_0^d\)，其中 \(|\alpha|=\alpha_1+\cdots+\alpha_d\)，\(\partial^\alpha=\partial_1^{\alpha_1}\cdots\partial_d^{\alpha_d}\) 在光滑函数上表示通常的偏导数。

\begin{definition}
\label{b1:def:weak-derivative}
给定 \(u,g\in L^1_{\mathrm{loc}}(\Omega)\) 和多指标 \(\alpha\)，若对每个 \(\varphi\in C_c^\infty(\Omega)\) 都有
\begin{equation}
\label{b1:eq:weak-derivative-test}
 \int_\Omega g\varphi\,\mathrm dx
 =(-1)^{|\alpha|}\int_\Omega u\,\partial^\alpha\varphi\,\mathrm dx,
\end{equation}
则称 \(g\) 为 \(u\) 的 \(\alpha\) 阶\term[ruo-daoshu]{弱导数}[weak derivative]，记为 \(\partial^\alpha u=g\)。当 \(\alpha=e_i\) 时，记为 \(\partial_i u\)。若全部一阶弱偏导数存在，便把它们合记为 \(\nabla u\coloneqq(\partial_1u,\ldots,\partial_du)\)。
\end{definition}
\symidx{\(\partial^\alpha u\)：局部可积函数的多指标弱导数}

所有积分只涉及测试函数的紧支集，所以局部可积性已经足够，不要求 \(u\) 或 \(g\) 在整个 \(\Omega\) 上可积。定义不含复共轭，与前面的复线性分布配对一致。当 \(|\alpha|=0\) 时，它只是函数本身；正阶时，定义并没有先要求经典差商存在。

用正则分布记号，\eqref{b1:eq:weak-derivative-test}就是 \(\partial^\alpha T_u=T_g\)。因此，本书把“\(u\) 有弱导数”限定为其分布导数能够由 \(L^1_{\mathrm{loc}}\) 函数表示；不能因为分布导数总存在，就认为弱导数也总存在。Hunter 与 Nachtergaele 的《Applied Analysis》\cite[Section 12.9, Definition 12.64]{Hunter2001Applied}采用这一局部可积函数的定义，并由此引入下一章的函数空间。

\subsection{唯一性}
\label{b1:subsec:200402}

\begin{proposition}
\label{b1:prop:weak-derivative-unique}
一个给定阶数的弱导数若存在，则在几乎处处相等的意义下唯一。把 \(u\) 或其弱导数在零测集上改值，不改变弱导数关系。
\end{proposition}
\begin{proof}
若 \(g_1,g_2\) 都满足定义，则 \(\int_\Omega(g_1-g_2)\varphi=0\) 对所有测试函数成立。由\cref{b1:prop:regular-distribution-injective}，\(g_1=g_2\) 几乎处处。零测集上的改值不影响定义中的积分，因此也不影响这个等价类。
\end{proof}

唯一性并不指定每一点的导数值。例如 \(u(x)=|x|\) 的一阶弱导数将在下文确定为 \(\operatorname{sgn}x\)，它在原点取哪个有限值都无关紧要。反过来，若要求导数具有额外的连续代表或指定的点值，就需要另外证明这种代表存在。

\begin{proposition}
\label{b1:prop:weak-derivative-locality}
设 \(u,g\in L^1_{\mathrm{loc}}(\Omega)\)，且 \(\{U_j\}\) 是 \(\Omega\) 的一个开覆盖。则 \(\partial^\alpha u=g\) 在 \(\Omega\) 上成立，当且仅当它在每个 \(U_j\) 上成立。
\end{proposition}
\begin{proof}
整体关系限制到开子集时仍成立，因为子集内的测试函数补零后仍是 \(\Omega\) 中的测试函数。反过来，分布 \(S=\partial^\alpha T_u-T_g\) 在每个 \(U_j\) 上为零。由\cref{b1:prop:distribution-locality}，\(S=0\) 在整个 \(\Omega\) 上成立，正是所需关系。
\end{proof}

弱导数关系还便于取极限。若 \(u_j\rightarrow u\)、\(g_j\rightarrow g\) 于 \(L^1_{\mathrm{loc}}(\Omega)\)，且 \(\partial^\alpha u_j=g_j\)，则 \(\partial^\alpha u=g\)。事实上，对固定测试函数，\(\varphi\) 与 \(\partial^\alpha\varphi\) 在同一个紧集上有界；两个积分的误差分别受相应的局部 \(L^1\) 误差乘这个界控制，因而可在\eqref{b1:eq:weak-derivative-test}中直接取极限。

这里需要同时控制函数和导数，而不是从 \(u_j\) 的收敛自动推断导数收敛。例如 \(u_j(x)=j^{-1}\sin(j^2x)\) 一致趋于零，但其经典导数 \(j\cos(j^2x)\) 在原点甚至趋于无穷。分布导数仍由前节的连续性趋于零；要求导数在某个函数空间中收敛，则是另外一个有实质内容的条件。

\subsection{分部积分}
\label{b1:subsec:200403}

\begin{proposition}[Leibniz]
\label{b1:prop:weak-leibniz}
若 \(u\) 的全部 \(|\beta|\leq r\) 阶弱导数局部可积，且 \(a\in C^\infty(\Omega)\)，则对 \(|\alpha|\leq r\)，
\[
 \partial^\alpha(au)
 =\sum_{\beta\leq\alpha}\binom\alpha\beta
    (\partial^{\alpha-\beta}a)(\partial^\beta u),
 \quad
 \binom\alpha\beta=\prod_{i=1}^d\binom{\alpha_i}{\beta_i}.
\]
其中 \(\beta\leq\alpha\) 指逐分量不等式，等式按弱导数解释。
\end{proposition}
\begin{proof}
先设 \(r=1\)，记 \(g_i=\partial_i u\)。将 \(a\varphi\) 代入定义，得到
\[
 \begin{aligned}
 -\int_\Omega au\,\partial_i\varphi
 &=-\int_\Omega u\,\partial_i(a\varphi)
    +\int_\Omega u(\partial_i a)\varphi\\
 &=\int_\Omega\bigl(ag_i+u\partial_i a\bigr)\varphi.
 \end{aligned}
\]
右侧函数局部可积，故这就证明一阶公式。高阶时，分布导数可交换，所以已存在的 \(\partial^\beta u\) 的弱偏导数为 \(\partial^{\beta+e_i}u\)。反复应用一阶公式；对每个坐标，\(\alpha_i\) 次求导中有 \(\beta_i\) 次落在 \(u\) 上的选择数为 \(\binom{\alpha_i}{\beta_i}\)，各坐标的选择数相乘，便得到所列系数及公式。
\end{proof}

上述计算是内部的分部积分，不要求 \(\Omega\) 的边界光滑。它也允许把测试函数类扩大到 \(C_c^1(\Omega)\)：若 \(\partial_i u=g_i\in L^1_{\mathrm{loc}}\)，则
\begin{equation}
\label{b1:eq:weak-integration-by-parts}
 \int_\Omega u\,\partial_i v\,\mathrm dx
 =-\int_\Omega g_i v\,\mathrm dx
 \quad(v\in C_c^1(\Omega)).
\end{equation}
为证明这一点，把 \(v\) 在域外补零，并用\secref{b1:sec:2003}的磨光函数取 \(v_\varepsilon=v*\rho_\varepsilon\)。当 \(\varepsilon\) 足够小时，这些函数的支集包含在同一个 \(K\Subset\Omega\) 中，且 \(v_\varepsilon\rightarrow v\)、\(\partial_i v_\varepsilon\rightarrow\partial_i v\) 一致收敛。这两个一致收敛也可直接由 \(v\) 及其一阶导数的一致连续性和 \(\int\rho_\varepsilon=1\) 得到。对 \(v_\varepsilon\) 使用定义，再由 \(u,g_i\in L^1(K)\) 取极限，得到\eqref{b1:eq:weak-integration-by-parts}。

紧支集条件在这里有实际作用：测试函数在边界附近为零，所以不出现边界项。对一直延伸到边界的函数使用分部积分，需要边界值及迹的概念，不能从上式省略紧支集条件后直接得到。

与此相应，在开子集上限制弱导数是安全的，在域外补零却需要额外检查。例如常函数 \(u=1\) 在 \((0,1)\) 上的弱导数为零，但其零延拓 \(\overline u=\mathbf1_{(0,1)}\) 满足
\[
 \langle\partial T_{\overline u},\varphi\rangle
 =-\int_0^1\varphi'(x)\,\mathrm dx
 =\varphi(0)-\varphi(1).
\]
因此零延拓产生了 \(\delta_0-\delta_1\) 两个端点项，并不是把内部的零导数也补零。若先乘以 \(\chi\in C_c^\infty(\Omega)\)，则 \(\chi u\) 在边界附近已经为零，Leibniz 公式给出的导数便可以与函数一起补零。这是局部化与直接截断定义域的区别。

\subsection{经典导数与局部 Lipschitz 函数}
\label{b1:subsec:200404}

若 \(u\in C^1(\Omega)\)，则经典偏导数也是弱偏导数。事实上，\(u\varphi\) 在域外补零后属于 \(C_c^1(\mathbb R^d)\)。沿第 \(i\) 个坐标积分，由一维微积分基本定理和 Fubini 定理，\(\int\partial_i(u\varphi)=0\)，展开乘积即得定义中的等式。重复这个论证，也能处理 \(C^r\) 函数的 \(r\) 阶导数。

\begin{theorem}
\label{b1:thm:lipschitz-weak-derivative}
局部 Lipschitz 函数 \(u:\Omega\rightarrow\mathbb R\) 或 \(\mathbb C\) 的各一阶弱导数存在，属于 \(L^\infty_{\mathrm{loc}}(\Omega)\)，并几乎处处等于经典偏导数。在 \(u\) 为 \(L\)-Lipschitz 的开集上，\(|\partial_i u|\leq L\) 几乎处处。
\end{theorem}
\begin{proof}
先在 \(u\) 为 \(L\)-Lipschitz 的有界开球 \(B\Subset\Omega\) 内证明。由\cref{b1:thm:rademacher}，经典微分几乎处处存在；再由\cref{b1:prop:lipschitz-differentiability-borel}，把不可微点处的微分矩阵补零后得到 Borel 矩阵函数。令 \(g_i\) 为其第 \(i\) 个偏导分量，则 \(g_i\) 可测且 \(|g_i|\leq L\)。将 \(u|_B\) 在球外补零为 \(u_0\)。对 \(\varphi\in C_c^\infty(B)\)，平移变量给出
\[
 \int_{\mathbb R^d}\dfrac{u_0(x+he_i)-u_0(x)}h\varphi(x)\,\mathrm dx
 =\int_{\mathbb R^d}u_0(x)
   \dfrac{\varphi(x-he_i)-\varphi(x)}h\,\mathrm dx.
\]
当 \(|h|\) 足够小时，左侧非零测试点所涉及的两点都在球内，差商被 \(L\) 控制，并几乎处处趋于 \(g_i\)。右侧测试差商支集留在固定的紧子集内，被 \(\|\partial_i\varphi\|_\infty\) 控制，且趋于 \(-\partial_i\varphi\)。两边应用支配收敛定理，得 \(\int_Bg_i\varphi=-\int_Bu\partial_i\varphi\)。最后用这样的球覆盖 \(\Omega\)，由弱导数的局部性和唯一性拼合。
\end{proof}

因此 \(|x|\) 的弱导数为 \(\operatorname{sgn}x\)，尖点不妨碍一阶弱导数存在。一维情形还有一个更精确的刻画，它同时说明绝对连续性为什么会出现。

\subsection{一维弱导数与绝对连续代表}
\label{b1:subsec:200405}

\begin{theorem}[一维弱导数与绝对连续代表]
\label{b1:thm:one-dimensional-weak-ac}
设 \(I\subseteq\mathbb R\) 为非空开区间，\(u,g\in L^1_{\mathrm{loc}}(I)\)。则 \(u'=g\) 作为弱导数成立，当且仅当 \(u\) 有一个代表 \(\widetilde u\)，在每个紧子区间上绝对连续，且 \(\widetilde u'=g\) 几乎处处。这个连续代表唯一，并满足
\[
 \widetilde u(t)-\widetilde u(s)=\int_s^t g(r)\,\mathrm dr
 \quad(s,t\in I).
\]
\end{theorem}
\begin{proof}
先证实值情形。复值情形分别对 \(u,g\) 的实部和虚部应用实值结论，再把所得两个绝对连续代表组合起来；因而以下构造中的函数均取实值。

先证明：若 \(w\in L^1_{\mathrm{loc}}(I)\) 的弱导数为零，则 \(w\) 几乎处处为常数。对积分为零的 \(\psi\in C_c^\infty(I)\)，先在 \(I\) 外补零，其原函数 \(\Psi(x)=\int_{-\infty}^x\psi(t)\,\mathrm dt\) 仍属于 \(C_c^\infty(I)\)，故 \(\int_Iw\psi=\int_Iw\Psi'=0\)。取固定的 \(\eta\in C_c^\infty(I)\)，使 \(\int_I\eta=1\)。对任意测试函数 \(\varphi\)，将上述结论用于 \(\psi=\varphi-(\int_I\varphi)\eta\)，便得
\[
 \int_I w\varphi=c\int_I\varphi,\quad c=\int_Iw\eta.
\]
正则分布嵌入的单射性说明 \(w=c\) 几乎处处。

现固定 \(x_0\in I\)，定义 \(G(x)=\int_{x_0}^xg(t)\,\mathrm dt\)。由\cref{b1:thm:ac-fundamental-calculus}，\(G\) 在每个紧子区间上绝对连续，且 \(G'=g\) 几乎处处。它也有弱导数 \(g\)：给定测试函数，取包含其支集的 \([a,b]\subset I\)，写 \(G(x)=G(a)+\int_a^xg(t)\,\mathrm dt\)，由 Fubini 定理，
\[
 \int_a^bG(x)\varphi'(x)\,\mathrm dx
 =\int_a^b g(t)\left(\int_t^b\varphi'(x)\,\mathrm dx\right)\mathrm dt
 =-\int_a^b g(t)\varphi(t)\,\mathrm dt.
\]
若 \(u'=g\)，则 \(u-G\) 的弱导数为零，所以 \(u=c+G\) 几乎处处，所需代表为 \(\widetilde u=c+G\)。反向则由绝对连续函数的积分表示及同一个 Fubini 计算得到。两个连续代表若几乎处处相等，则处处相等，否则连续性会给出一个它们不相等的开子区间，故代表唯一。
\end{proof}

这个刻画给出了超出 Lipschitz 类的具体函数。若 \(0<\beta<1\)，则 \(u(x)=|x|^\beta\) 在原点附近不是 Lipschitz 函数，但
\[
 |x|^\beta=\int_0^x\beta\operatorname{sgn}(t)|t|^{\beta-1}\,\mathrm dt.
\]
被积函数局部可积，因而它就是 \(u\) 的弱导数。进一步，对 \(1\leq p<\infty\)，该导数在原点附近属于 \(L^p\)，当且仅当 \(p(1-\beta)<1\)；这可由积分 \(\int_0^r t^{-p(1-\beta)}\,\mathrm dt\) 直接判断。在临界等号处，积分按对数发散。导数存在与导数具有何种可积性，是两个不同层次的问题。

最后看两个反例。令 \(H=\mathbf1_{(0,\infty)}\)。它的经典导数在原点外为零，但
\[
 \langle\partial T_H,\varphi\rangle
 =-\int_0^\infty\varphi'(x)\,\mathrm dx
 =\varphi(0),
\]
所以分布导数是 \(\delta_0\)，而不是零。它不能由局部可积函数表示：若能，则相应函数在不含原点的开集上由正则嵌入的单射性为零，于是几乎处处为零，却无法对满足 \(\varphi(0)=1\) 的测试函数给出值 \(1\)。因此 \(H\) 没有本节意义下的一阶弱导数。

\cref{b1:fig:jump-distributional-derivative}把这一差别画成最简单的模型：经典导数在跳跃点外为零，分布导数却把全部跳跃质量集中到该点。
\begin{figure}[htbp]
  \centering
  \begin{tikzpicture}[>=stealth]
    \begin{scope}
      \draw[->] (-1.8,0)--(1.9,0);
      \draw[->] (0,-0.25)--(0,1.45);
      \draw[thick] (-1.55,0)--(0,0);
      \draw[thick] (0,1)--(1.55,1);
      \draw[dashed] (0,0)--(0,1);
      \fill[white] (0,0) circle (2pt);
      \draw (0,0) circle (2pt);
      \fill (0,1) circle (2pt);
      \node[above right] at (0,1) {$H$};
    \end{scope}
    \draw[->,thick] (2.25,0.62)--(3.65,0.62) node[midway,above] {$\partial$};
    \begin{scope}[shift={(5.6,0)}]
      \draw[->] (-1.5,0)--(1.55,0);
      \draw[->] (0,-0.25)--(0,1.45);
      \draw[->,very thick] (0,0)--(0,1.15);
      \node[above right] at (0,1.05) {$\delta_0$};
    \end{scope}
  \end{tikzpicture}
  \caption[跳跃的分布导数]{跳跃的分布导数。Heaviside 函数的变化集中在原点，故 \(\partial T_H=\delta_0\)。}
  \label{b1:fig:jump-distributional-derivative}
\end{figure}

连续性也不能排除这种问题。\cref{b1:ex:14-cantor-function}构造的 Cantor 函数 \(F_C\) 连续、非恒定，且经典导数几乎处处为零。如果它在 \((0,1)\) 上有局部可积的弱导数，由刚证的定理，便有局部绝对连续代表；该代表与连续的 \(F_C\) 必须处处相等。其弱导数于是等于经典导数，几乎处处为零，迫使 \(F_C\) 为常数，矛盾。跳跃函数和 Cantor 函数分别说明，点上的跳跃与零测集上的连续变化，都可能被经典的几乎处处导数遗漏。

% 本章小结（不计入编号）
\begin{chaptersummary}
紧支集光滑函数允许任意次分部积分，又把计算限制在开集内部。分布的连续性在每个固定紧集上化成一个有限阶估计，因此可以把测试函数上的微分、光滑乘法和局部化转移到分布上。一般分布无需有点值，也无需由可积密度给出。

光滑化使分布暂时成为普通光滑函数，并保留微分关系；其收敛首先是逐个测试函数的配对收敛，不能自动理解为一致收敛或逐点收敛。弱导数则进一步要求所得分布由局部可积函数表示。在一维，这与局部绝对连续代表相联系；跳跃产生点质量，Cantor 函数提醒我们不能只凭几乎处处的经典导数判断整个变化。下一章将对函数和弱导数同时施加 \(L^p\) 控制，建立 Sobolev 空间。
\end{chaptersummary}


\BookExercises

% \chapref{b1:ch:20}习题临时稿；不另设 BookExercises 入口。
% 十一题均为依据本章已建立工具自拟的变式，不冒认教材原题编号。
% 训练分工：测试拓扑/局部阶；乘法核/主值构造；跳跃/分布ODE；磨光误差/弱导数端点。
% 前置：20.1测试列判据，20.2有限阶与分布运算，20.3卷积，20.4弱导数；无21章依赖。

\begin{exercise}\label{b1:ex:20-basic-distribution-pairings}
在 \(\mathcal D'(\mathbb R)\) 中按定义计算下列对象对测试函数 \(\varphi\) 的作用，并辨认所得分布：
\[
 \delta_a',\qquad x\delta_0',\qquad (T_x)'.
\]
其中 \(T_x\) 表示函数 \(x\mapsto x\) 产生的正则分布。
\end{exercise}
% 来源：自拟；先以三项按定义配对训练点导数、光滑乘法与正则分布求导。
\begin{solution}
分布导数的定义直接给出
\[
 \langle\delta_a',\varphi\rangle=-\varphi'(a).
\]
再由光滑函数与分布的乘法，
\[
 \langle x\delta_0',\varphi\rangle
 =\langle\delta_0',x\varphi\rangle
 =-(x\varphi)'(0)=-\varphi(0),
\]
所以 \(x\delta_0'=-\delta_0\)。最后作一次普通分部积分，测试函数的边界项为零，故
\[
 \langle(T_x)',\varphi\rangle
 =-\int_{\mathbb R}x\varphi'(x)\,\mathrm dx
 =\int_{\mathbb R}\varphi(x)\,\mathrm dx
 =\langle T_1,\varphi\rangle.
\]
因此 \((T_x)'=T_1\)。
\end{solution}

\begin{exercise}\label{b1:ex:20-basic-weak-derivatives}
在 \(\mathbb R\) 上分别考虑 \(1\)、\(x\)、\(|x|\) 和 \(H=\mathbf1_{(0,\infty)}\)。
\begin{enumerate}
\item 求前三个函数的一阶弱导数。
\item 证明 \(H\) 的分布导数为 \(\delta_0\)，因而它没有局部可积的一阶弱导数。
\item 求 \(|x|\) 的二阶分布导数，并判断它是否存在局部可积的二阶弱导数。
\end{enumerate}
\end{exercise}
% 来源：自拟；用最基本模型区分经典不可微点、弱导数和集中型分布导数。
\begin{solution}
常函数与恒等函数的经典导数分别为 \(0\) 和 \(1\)，由经典导数与弱导数的一致性，它们的一阶弱导数也是这两个函数。函数 \(|x|\) 局部 Lipschitz，所以由\cref{b1:thm:lipschitz-weak-derivative}，
\[
 (|x|)'=\operatorname{sgn}x
 \quad\text{几乎处处}.
\]
原点处给 \(\operatorname{sgn}x\) 赋哪个有限值都不影响弱导数。

对任意测试函数，
\[
 \langle(T_H)',\varphi\rangle
 =-\int_0^\infty\varphi'(x)\,\mathrm dx
 =\varphi(0)=\langle\delta_0,\varphi\rangle.
\]
若 \(H\) 有局部可积弱导数，该函数产生的正则分布就应等于 \(\delta_0\)，与\cref{b1:prop:dirac-not-regular}矛盾。

同理，分段积分给出
\[
 \langle(\operatorname{sgn}x)',\varphi\rangle=2\varphi(0),
\]
所以 \((T_{|x|})''=2\delta_0\)。它不是正则分布，故 \(|x|\) 没有本章意义下的二阶局部可积弱导数。尖点容许一阶弱导数，却在再求一次导数时产生点质量。
\end{solution}

\begin{exercise}\label{b1:ex:20-dirac-mollification}
设 \(a\in\mathbb R^d\)，\(\rho_\varepsilon\) 为本章的磨光函数。证明
\[
 \delta_a*\rho_\varepsilon=\rho_\varepsilon(\,\cdot-a),
 \qquad
 (\partial^\alpha\delta_a)*\rho_\varepsilon
 =\partial^\alpha\rho_\varepsilon(\,\cdot-a).
\]
再直接验证这两个光滑函数产生的正则分布分别趋于 \(\delta_a\) 和 \(\partial^\alpha\delta_a\)。
\end{exercise}
% 来源：自拟；直接练习点质量的卷积、导数交换与分布收敛。
\begin{solution}
由分布卷积的定义，
\[
 (\delta_a*\rho_\varepsilon)(x)
 =\langle(\delta_a)_y,\rho_\varepsilon(x-y)\rangle
 =\rho_\varepsilon(x-a).
\]
卷积微分公式随即给出第二个等式。对任意 \(\varphi\in\mathcal D(\mathbb R^d)\)，变量代换 \(x=a+\varepsilon z\) 得
\[
 \int_{\mathbb R^d}\rho_\varepsilon(x-a)\varphi(x)\,\mathrm dx
 =\int_{\mathbb R^d}\rho(z)\varphi(a+\varepsilon z)\,\mathrm dz
 \longrightarrow\varphi(a),
\]
其中使用 \(\rho\) 的紧支集、单位积分和 \(\varphi\) 的连续性。因此相应正则分布趋于 \(\delta_a\)。对导数再作分部积分，
\[
 \int\partial^\alpha\rho_\varepsilon(x-a)\varphi(x)\,\mathrm dx
 =(-1)^{|\alpha|}\int\rho_\varepsilon(x-a)\partial^\alpha\varphi(x)\,\mathrm dx
 \longrightarrow(-1)^{|\alpha|}\partial^\alpha\varphi(a),
\]
右端正是 \(\langle\partial^\alpha\delta_a,\varphi\rangle\)。
\end{solution}

\begin{exercise}\label{b1:ex:20-test-convergence-boundaries}
\begin{enumerate}
\item\label{b1:item:20-escaping-flat-amplitude}
取 \(0\ne\psi\in C_c^\infty\bigl((-1/4,1/4)\bigr)\)，在 \(\Omega=(0,1)\) 上令
\[
 \varphi_n(x)=e^{-n^2}\psi\bigl(n^2(x-1/n)\bigr),
 \quad n\ge4.
\]
证明它的每个固定阶导数在整个 \(\Omega\) 上一致趋零，但 \(\varphi_n\) 不在 \(\mathcal D(\Omega)\) 中趋于零。
\item\label{b1:item:20-oscillatory-regular-distributions}
取 \(\chi\in\mathcal D\bigl((-2,2)\bigr)\)，使 \(\chi=1\) 于 \([-1,1]\)，令 \(u_n(x)=n^{-1}\chi(x)\sin(nx)\)。证明 \(T_{u_n}\to0\) 于 \(\mathcal D'\bigl((-2,2)\bigr)\)，但 \(u_n\not\to0\) 于测试函数空间。
\end{enumerate}
\end{exercise}
% 来源：自拟；逃逸例增加超快衰减，使所有全域导数上确界都趋零；高频例比较测试对象与正则分布两种收敛。
\begin{solution}
第一问中，\(\varphi_n\) 的支集落在 \([1/n-1/(4n^2),1/n+1/(4n^2)]\subset(0,1)\)。对每个固定整数 \(j\ge0\)，
\[
 \|\varphi_n^{(j)}\|_\infty
 =e^{-n^2}n^{2j}\|\psi^{(j)}\|_\infty\longrightarrow0.
\]
然而取 \(z_0\) 使 \(\psi(z_0)\ne0\)，则 \(x_n=1/n+z_0/n^2\) 为非零取值点，且 \(x_n\to0\)。而 \((0,1)\) 的任一紧子集都与端点保持正距离，所以不存在共同紧集包含全部支集。由\cref{b1:prop:test-sequence-convergence}，测试函数收敛不成立。振幅再小，也不能消除支集逃向边界的问题。

第二问的全部支集都在 \(\operatorname{supp}\chi\) 中。对任意固定测试函数 \(\eta\)，
\[
 |\langle T_{u_n},\eta\rangle|
 \le\frac{\|\chi\|_\infty}{n}\int_{-2}^2|\eta(x)|\,\mathrm dx\longrightarrow0.
\]
但在 \([-1,1]\) 上有 \(u_n'(x)=\cos(nx)\)，特别是 \(u_n'(0)=1\)。因此一阶导数不一致趋零，仍不满足测试列判据。分布收敛只逐个测试其积分，不能据此反推函数及其导数的测试空间收敛。
\end{solution}

\begin{optionalexercise}\label{b1:ex:20-locally-finite-unbounded-order}
在 \(\mathbb R\) 上定义线性泛函
\[
 \langle T,\varphi\rangle
 =\sum_{n=1}^\infty(-1)^n\varphi^{(n)}(n),
 \quad \varphi\in\mathcal D(\mathbb R).
\]
证明它是支集恰为 \(\{1,2,\ldots\}\) 的分布，但不存在一个固定整数 \(m\)，使其在每个紧集上都满足阶数不超过 \(m\) 的估计。这里允许估计常数随紧集改变。
\end{optionalexercise}
% 来源：自拟；将局部有限的点支集和逐渐增加的导数阶数组合，具体实现正文关于局部阶的警告。
\begin{solution}
测试函数的紧支集只包含有限个正整数，其余点处的函数及全部导数为零，所以这个和实际上有限。对固定紧集 \(K\)，若其中最大的正整数为 \(N\)，则
\[
 |\langle T,\varphi\rangle|\le Np_{K,N}(\varphi)
 \quad(\varphi\in\mathcal D_K).
\]
若 \(K\) 不含正整数，左端恒为零。局部有限阶估计证明 \(T\in\mathcal D'(\mathbb R)\)。

在不含正整数的开集上，\(T\) 为零；在正整数 \(N\) 附近，它与 \(\delta_N^{(N)}\) 相同。取支集足够小、在 \(N\) 处的 \(N\) 阶导数非零的测试函数，便知它在 \(N\) 的任一邻域不恒为零。因此支集恰为所列离散集合。

假设存在全局通用阶数 \(m\)。选整数 \(N>m\)，固定 \(K=[N-1/4,N+1/4]\)，再取 \(\psi\in C_c^\infty\bigl((-1/4,1/4)\bigr)\) 使 \(\psi^{(N)}(0)=1\)。对 \(0<\varepsilon<1\)，令
\[
 \psi_\varepsilon(x)=\varepsilon^m\psi\left(\frac{x-N}{\varepsilon}\right).
\]
其支集在 \(K\) 中，且所有 \(j\le m\) 阶导数的上确界均有独立于 \(\varepsilon\) 的界。但
\[
 |\langle T,\psi_\varepsilon\rangle|=\varepsilon^{m-N}\longrightarrow\infty,
\]
与 \(|\langle T,\varphi\rangle|\le C_Kp_{K,m}(\varphi)\) 矛盾。这不是常数在远处增大的问题，而是在同一个选定紧集上，第 \(m\) 阶控制已经不够。
\end{solution}

\begin{optionalexercise}\label{b1:ex:20-polynomial-annihilator}
设整数 \(m\ge1\)，\(T\in\mathcal D'(\mathbb R)\)。证明
\[
 x^mT=0
 \quad\Longleftrightarrow\quad
 T=\sum_{j=0}^{m-1}c_j\delta_0^{(j)}
\]
其中系数 \(c_j\in\mathbb C\) 唯一。解答须说明所用的除以 \(x^m\) 的商为何光滑且有紧支集，不预借单点支集分布的一般结构定理。并求 \(x\delta_0^{(j)}\)。
\end{optionalexercise}
% 来源：自拟；用有限阶 Taylor 分解求光滑乘法算子的核，实质扩展正文的一阶点质量恒等式。
\begin{solution}
固定 \(\chi\in\mathcal D(\mathbb R)\)，使 \(\chi=1\) 于原点附近。对任意测试函数 \(\varphi\)，令
\[
 r(x)=\varphi(x)-\chi(x)\sum_{j=0}^{m-1}\frac{\varphi^{(j)}(0)}{j!}x^j.
\]
函数 \(r\) 有紧支集，且原点处低于 \(m\) 阶的导数全部为零。在原点附近，带积分余项的 Taylor 公式给出
\[
 \frac{r(x)}{x^m}
 =\frac1{(m-1)!}\int_0^1(1-t)^{m-1}\varphi^{(m)}(tx)\,\mathrm dt.
\]
右端在 \(x=0\) 也光滑，任意阶微分可在紧区间积分下进行。因此这个商在原点光滑延拓；在原点外按通常除法定义，且因 \(r\) 紧支集而仍有紧支集。记所得测试函数为 \(\psi\)。

若 \(x^mT=0\)，则 \(\langle T,x^m\psi\rangle=0\)，所以
\[
 \langle T,\varphi\rangle
 =\sum_{j=0}^{m-1}\frac{\langle T,x^j\chi\rangle}{j!}\varphi^{(j)}(0).
\]
与 \(\langle\delta_0^{(j)},\varphi\rangle=(-1)^j\varphi^{(j)}(0)\) 比较，即得
\[
 c_j=\frac{(-1)^j}{j!}\langle T,x^j\chi\rangle.
\]
函数 \(x^j\chi/j!\) 在原点的各阶导数可分别取出相应系数，故系数唯一。

对于 \(j\ge1\)，乘积求导给出
\[
 \langle x\delta_0^{(j)},\varphi\rangle
 =(-1)^j(x\varphi)^{(j)}(0)
 =(-1)^jj\varphi^{(j-1)}(0),
\]
所以 \(x\delta_0^{(j)}=-j\delta_0^{(j-1)}\)，而 \(x\delta_0=0\)。迭代可知 \(x^m\delta_0^{(j)}=0\) 对 \(j<m\) 成立，这证明了反向蕴涵。
\end{solution}

\begin{optionalexercise}\label{b1:ex:20-principal-value-logarithm}
对 \(\varphi\in\mathcal D(\mathbb R)\)，考察
\[
 \langle P,\varphi\rangle
 =\lim_{\varepsilon\downarrow0}
 \int_{|x|>\varepsilon}\frac{\varphi(x)}x\,\mathrm dx.
\]
\begin{enumerate}
\item\label{b1:item:20-principal-value-existence}
证明极限存在并定义一个分布，且 \(P=(T_{\log|x|})'\)。
\item\label{b1:item:20-principal-value-nonregular}
证明 \(P\) 不是正则分布，\(xP=T_1\)，并求方程 \(xT=T_1\) 的全部分布解。
\end{enumerate}
\end{optionalexercise}
% 来源：自拟；对称截断构造及其边界项均在此独立证明，不预借 Hilbert 变换。
\begin{solution}
固定 \(R>0\)，使 \(\operatorname{supp}\varphi\subset(-R,R)\)。对称截断可写为
\[
 \int_{\varepsilon<|x|<R}\frac{\varphi(x)}x\,\mathrm dx
 =\int_\varepsilon^R\frac{\varphi(x)-\varphi(-x)}x\,\mathrm dx.
\]
由微积分基本定理，分子绝对值不超过 \(2x\|\varphi'\|_\infty\)，故极限存在，而且
\[
 |\langle P,\varphi\rangle|\le2R\|\varphi'\|_\infty.
\]
对每个固定紧支集可统一选 \(R\)，这证明分布的连续性。

函数 \(\log|x|\) 在零点附近可积。在 \((-R,-\varepsilon)\) 与 \((\varepsilon,R)\) 上分别分部积分，得到
\[
 -\int_{\varepsilon<|x|<R}\log|x|\,\varphi'(x)\,\mathrm dx
 =\int_{\varepsilon<|x|<R}\frac{\varphi(x)}x\,\mathrm dx
 +\log\varepsilon\bigl(\varphi(\varepsilon)-\varphi(-\varepsilon)\bigr).
\]
最后一项的绝对值至多为 \(2\varepsilon|\log\varepsilon|\cdot\|\varphi'\|_\infty\)，趋于零；左端因局部可积性趋于完整积分。这就证明 \((T_{\log|x|})'=P\)。

若 \(P=T_g\) 且 \(g\in L^1_{\mathrm{loc}}\)，则在不含原点的开区间上，测试积分给出 \(g=1/x\) 几乎处处。用可数个这样的区间覆盖 \(\mathbb R\setminus\{0\}\)，便得 \(g=1/x\) 几乎处处，与零点附近的绝对可积性矛盾。因此 \(P\) 非正则。

直接测试可得
\[
 \langle xP,\varphi\rangle
 =\lim_{\varepsilon\downarrow0}\int_{|x|>\varepsilon}\varphi(x)\,\mathrm dx
 =\langle T_1,\varphi\rangle.
\]
所以任意解与 \(P\) 的差满足 \(x(T-P)=0\)。由\cref{b1:ex:20-polynomial-annihilator}，全部解恰为 \(P+c\delta_0\)，其中 \(c\in\mathbb C\)。
\end{solution}

\begin{exercise}\label{b1:ex:20-piecewise-jump-derivatives}
设 \(a_1<\cdots<a_N\)，函数 \(f\) 在这些点分出的各开区间内为 \(C^2\) 函数，且在每个有限端点有有限的 \(f,f',f''\) 单侧连续延拓。在断点任意赋值。用 \(f'_{\mathrm{cl}}\)、\(f''_{\mathrm{cl}}\) 表示各段经典导数，并记
\[
 [f]_{a_j}=f(a_j+)-f(a_j-),
 \quad [f']_{a_j}=f'(a_j+)-f'(a_j-).
\]
证明
\[
 (T_f)'=T_{f'_{\mathrm{cl}}}+\sum_{j=1}^N[f]_{a_j}\delta_{a_j},
\]
\[
 (T_f)''=T_{f''_{\mathrm{cl}}}
 +\sum_{j=1}^N[f']_{a_j}\delta_{a_j}
 +\sum_{j=1}^N[f]_{a_j}\delta_{a_j}'.
\]
据此给出 \(f\) 存在一阶、二阶局部可积弱导数的充要条件。
\end{exercise}
% 来源：自拟；将单跳跃检测扩展为有限多界面并区分函数跳跃和一阶导数跳跃的分布阶数。
\begin{solution}
对固定测试函数，在其支集外选左右端点，再在所有 \(a_j\) 处分段。在一个相邻断点之间的区间 \((a,b)\) 上，分部积分给出
\[
 -\int_a^bf\varphi'
 =\int_a^bf'_{\mathrm{cl}}\varphi
 +f(a+)\varphi(a)-f(b-)\varphi(b).
\]
求和后，外端点项因测试函数在那里为零而消失；每个内断点的两个贡献合成 \([f]_{a_j}\varphi(a_j)\)，所以一阶公式成立。将相同计算用于分段函数 \(f'_{\mathrm{cl}}\)，并对点质量项求导，就得到二阶公式。

各段的经典导数均局部可积。有限点支集的正则分布只能为零：在这些点的补集上它为零，正则嵌入单射性迫使其代表几乎处处为零。另一方面，可在每个断点附近独立选择测试函数的值与一阶导数，从而分别取出 \(\delta_{a_j}\) 与 \(\delta_{a_j}'\) 的系数。因此一阶分布导数为正则分布，当且仅当所有 \([f]_{a_j}=0\)；二阶分布导数为正则分布，当且仅当所有 \([f]_{a_j}\) 和 \([f']_{a_j}\) 都为零。

故一阶弱导数存在等价于各段函数值能连续拼接；二阶弱导数存在等价于函数值及一阶导数都能连续拼接。断点上预先填写的单点值不改变这些结论，因为弱导数只依赖几乎处处等价类。
\end{solution}

\begin{optionalexercise}\label{b1:ex:20-distribution-first-order-ode}
设 \(I\subset\mathbb R\) 为非空开区间，\(a\in\mathbb C\)。证明 \(\mathcal D'(I)\) 中的方程
\[
 T'+aT=0
\]
的每个解都是光滑函数 \(ce^{-ax}\) 所产生的正则分布。不得预先假设 \(T\) 有函数代表。若 \(I\) 改成不连通开集，常数应如何选择？
\end{optionalexercise}
% 来源：自拟；把积分因子方法推广到一般分布，补证零导数分布为常数，而非预借正则代表。
\begin{solution}
先设 \(S\in\mathcal D'(I)\) 且 \(S'=0\)。若 \(\psi\in\mathcal D(I)\) 且 \(\int_I\psi=0\)，将其补零后令
\[
 \Psi(x)=\int_{-\infty}^x\psi(t)\,\mathrm dt.
\]
在 \(\operatorname{supp}\psi\) 左右两侧它都为零，且其支集位于 \(I\) 内的紧区间，所以 \(\Psi\in\mathcal D(I)\)。于是
\[
 \langle S,\psi\rangle=\langle S,\Psi'\rangle=-\langle S',\Psi\rangle=0.
\]
取固定 \(\eta\in\mathcal D(I)\)，使 \(\int\eta=1\)。对任意 \(\varphi\)，函数 \(\varphi-(\int\varphi)\eta\) 积分为零，故
\[
 \langle S,\varphi\rangle=c\int_I\varphi,
 \quad c=\langle S,\eta\rangle.
\]
这证明一般分布 \(S\) 已是常函数 \(c\) 的正则分布。

对原方程令 \(S=e^{ax}T\)。由分布的光滑乘法及 Leibniz 公式，
\[
 S'=e^{ax}(T'+aT)=0.
\]
所以 \(e^{ax}T=T_c\)，乘以处处非零的 \(e^{-ax}\) 得 \(T=T_{ce^{-ax}}\)。反过来，这些光滑函数的经典导数满足原方程，因而也是分布解。

若底空间为不连通开集，每个连通分支都是开区间，可在各分支上独立选择一个常数。这些解能拼成局部光滑函数：一个紧子集可被有限个包含于开集内的区间覆盖，只会遇到有限个分支，故不存在分支端点处的内部跳跃条件。
\end{solution}

\begin{optionalexercise}\label{b1:ex:20-smoothing-variable-coefficient}
设 \(\Omega\subset\mathbb R^d\) 为开集，\(U\) 为开集且 \(K\Subset U\)、\(\overline U\Subset\Omega\)，\(a\in C^\infty(\Omega)\)，\(f\in L^p_{\mathrm{loc}}(\Omega)\)，\(1\le p\le\infty\)。沿用本章非负单位积分光滑核 \(\rho_\varepsilon\)，令
\[
 C_\varepsilon(x)=a(x)(f*\rho_\varepsilon)(x)-\bigl((af)*\rho_\varepsilon\bigr)(x).
\]
证明当 \(\varepsilon\) 充分小时，
\[
 \|C_\varepsilon\|_{L^p(K)}
 \le\varepsilon\sup_U|\nabla a|\cdot\|f\|_{L^p(U)}
       \int_{\mathbb R^d}|z|\rho(z)\,\mathrm dz.
\]
最后在一维取 \(T=\delta_0\)、\(a(x)=x\)，证明相应表达式
\[
 x(T*\rho_\varepsilon)-(xT)*\rho_\varepsilon
\]
虽在分布意义下趋于零，却不在原点的任一固定邻域上依 \(L^\infty\) 范数趋于零。
\end{optionalexercise}
% 来源：自拟；量化光滑乘法与卷积不交换，有限p和无穷p只用于误差估计，不宣称一般L∞磨光范数收敛。
\begin{solution}
选择 \(\varepsilon<\operatorname{dist}(K,U^c)\)。若 \(x\in K\)、\(|z|\le\varepsilon\)，则连接 \(x\) 与 \(x-z\) 的整条线段都在 \(U\) 中。沿线段积分给出
\[
 |a(x)-a(x-z)|\le |z|\sup_U|\nabla a|.
\]
把 \(f\mathbf1_U\) 在外部补零，记作 \(f_0\)。卷积的差可写为
\[
 C_\varepsilon(x)=\int\bigl(a(x)-a(x-z)\bigr)f_0(x-z)\rho_\varepsilon(z)\,\mathrm dz.
\]
因 \(\rho\ge0\)，其绝对值受 \(\sup_U|\nabla a|\) 乘以 \(k_\varepsilon*|f_0|\) 控制，其中 \(k_\varepsilon(z)=|z|\rho_\varepsilon(z)\)。由\cref{b1:lem:convolution-lp-bound}，
\[
 \|C_\varepsilon\|_{L^p(K)}
 \le\sup_U|\nabla a|\cdot\|k_\varepsilon\|_1\cdot\|f_0\|_p.
\]
而 \(\|k_\varepsilon\|_1=\varepsilon\int|z|\rho(z)\,\mathrm dz\)，这就给出所求估计，包括 \(p=\infty\)。

对点质量，\(x\delta_0=0\)、\(\delta_0*\rho_\varepsilon=\rho_\varepsilon\)，所以差等于 \(x\rho_\varepsilon(x)\)。在一维令 \(x=\varepsilon z\)，可得
\[
 \sup_x|x\rho_\varepsilon(x)|=\sup_z|z|\rho(z)>0.
\]
核非负且积分为 \(1\)，不能只集中在单点，故右端确实为正。支集随 \(\varepsilon\downarrow0\) 缩进任一固定原点邻域，因而局部 \(L^\infty\) 收敛失败。另一方面，
\[
 \|x\rho_\varepsilon(x)\|_{L^1(\mathbb R)}
 =\varepsilon\int|z|\rho(z)\,\mathrm dz\longrightarrow0,
\]
所以它在分布意义下趋于零。函数误差的范数结论不能直接推广给没有局部可积代表的一般分布。
\end{solution}

\begin{exercise}\label{b1:ex:20-power-weak-derivative-thresholds}
对实数 \(a>-1\)，在 \((-1,1)\) 上令 \(u_a(x)=|x|^a\)，原点处任取有限值；\(a=0\) 时按几乎处处恒等于 \(1\) 解释。确定哪些 \(a\) 使它分别有一阶、二阶局部可积弱导数，并在存在时给出公式。特别说明 \(a=1\) 时二阶分布导数是什么。
\end{exercise}
% 来源：自拟；把分段分部积分扩展到幂次奇性，端点a=1单独处理，不将经典导数代入零点来猜弱导数。
\begin{solution}
条件 \(a>-1\) 保证 \(u_a\) 在零点附近可积。若 \(a>0\)，候选一阶弱导数为
\[
 g_a(x)=a\operatorname{sgn}(x)|x|^{a-1}\quad(x\ne0),
\]
它也局部可积。在原点两侧截去 \((-\varepsilon,\varepsilon)\) 后分部积分；测试函数在外端点为零，所得内边界项是
\[
 \varepsilon^a\bigl(\varphi(\varepsilon)-\varphi(-\varepsilon)\bigr),
\]
其绝对值不超过 \(2\varepsilon^{a+1}\|\varphi'\|_\infty\)，趋于零。两侧积分由绝对可积性传到极限，证明 \(u_a'=g_a\) 为弱导数。当 \(a=0\) 时弱导数为零。

若 \(-1<a<0\) 且有局部可积弱导数 \(g\)，则在原点外的各开区间上，经典与弱导数一致性及唯一性迫使 \(g=g_a\) 几乎处处。但 \(\int_0^r|g_a|=\infty\)，与 \(g\in L^1_{\mathrm{loc}}\) 矛盾。因此一阶弱导数存在的充要条件是 \(a\ge0\)。

若 \(a>1\)，对 \(g_a\) 重复两侧分部积分，内边界项为
\[
 a\varepsilon^{a-1}\bigl(\varphi(\varepsilon)+\varphi(-\varepsilon)\bigr)\longrightarrow0.
\]
其经典导数 \(a(a-1)|x|^{a-2}\) 局部可积，所以
\[
 u_a''=a(a-1)|x|^{a-2}
\]
为二阶弱导数。当 \(a=0\) 时二阶导数仍为零。

若 \(a\ne0,1\) 且 \(-1<a\le1\)，任何局部可积的二阶弱导数在原点外都必须等于 \(a(a-1)|x|^{a-2}\)，但该函数在零点附近不可积，所以不存在。剩下 \(a=1\)：一阶弱导数为 \(\operatorname{sgn}x\)，上式中的边界项此时趋于 \(2\varphi(0)\)，两侧经典二阶导数为零，故
\[
 (T_{|x|})''=2\delta_0.
\]
它不是正则分布，因此 \(|x|\) 没有二阶局部可积弱导数。综上，二阶弱导数存在当且仅当 \(a=0\) 或 \(a>1\)。一阶导数在尖点处的跳跃，恰好成为二阶导数中的点质量。
\end{solution}
